%<*driver>
\def\nameofplainTeX{plain}
\ifx\fmtname\nameofplainTeX
\input l3docstrip.tex
\preamble

========================================================================
 ltxtalk-theme - themes for ltxtalk

 (C) 2026 Matthias Werner
 E-mail: matthias.werner@informatik.tu-chemnitz.de

 Released under the LaTeX Project Public License v1.3c or later
 See http://www.latex-project.org/lppl.txt
======================================================================== 

\endpreamble
\postamble

 This work is "maintained" (as per LPPL maintenance status) by
 Matthias Werner.
\endpostamble

\generate{
  \file{ltxtalk-theme.sty}{\from{lttheme.dtx}{pkg}}
  \file{ltxtalk-theme-minimal.sty}{\from{lttheme.dtx}{minimal}}
  \file{ltxtalk-theme-magpie.sty}{\from{lttheme.dtx}{magpie}}
  \file{ltxtalk-theme-hawk.sty}{\from{lttheme.dtx}{hawk}}
  \file{ltxtalk-theme-sparrow.sty}{\from{lttheme.dtx}{sparrow}}
  \file{ltxtalk-theme-anchovy.sty}{\from{lttheme.dtx}{anchovy}}
  \file{ltxtalk-theme-carp.sty}{\from{lttheme.dtx}{carp}}
  }
  \expandafter\endbatchfile
\fi
 \documentclass[
load-preamble+
]{osgdoc}

\usepackage{enumitem}
\usepackage{microtype}
\usepackage[
languages={de,en},
map={de=german,en=english},
targetlang={job=2},
load babel
]{langselect}

\def\pkgname{ltxtalk-theme}
\def\thepkg{\pkg{\pkgname}}
\newpackagename{\lttheme}{\pkgname}
\newpackagename{\lttalk}{ltx-talk}
\newpackagename{\beamer}{beamer}

\makeindex
\setcnltx{
  name  =  \pkgname,
  version  = \UseVersionOf{ltxtalk-theme.sty},
  date     = \UseDateOf{ltxtalk-theme.sty},
  title    = {\ldeen{Das \lttheme\ Paket}{The \lttheme\ Package}},
  info     = {\ldeen{Themen-}{Theme }Engine \ldeen{für}{for} \lttalk},
  authors  = {Matthias Werner[matthias.werner@informatik.tu-chemnitz.de]},
  abstract={
  \ldeen{   
      Die Klasse \lttalk{} ist eine \LaTeX-Klasse für Präsentationen und soll 
      mittelfristig die verbreitete Klasse \beamer{} ablösen, um PDF-Tagging 
      für Präsentationen zu ermöglichen.
      Das Package \lttheme\ unterstützt Entwickler, die Themen für \lttalk{} schreiben.
    }{
      The \lttalk{} class is a \LaTeX class for presentations and is intended 
      to eventually replace the widely used \beamer{} class in order to enable PDF tagging 
      for presentations.
      The \lttheme\ package supports developers who write themes for \lttalk.  
    }
  },
  url      =https://github.com/werner-matthias/osglecture,
  build-title
}
\begin{document}
% --------------------------------------------------------------------------
% Documentation
% --------------------------------------------------------------------------

\section{\ldeen{Einführung}{Introduction}}

\subsection{\ldeen{Ziel des Pakets}{Purpose of the package}}

\ldeen*{
  Die Klasse @1 ist eine \LaTeX-Klasse für Präsentationen und soll 
  mittelfristig die verbreitete Klasse @2 ablösen, um PDF-Tagging 
  für Präsentationen zu ermöglichen.
}{
  The @1 class is a \LaTeX class for presentations and is intended 
  to eventually replace the widely used @2 class in order to enable PDF tagging 
  for presentations.
}{\lttalk\footnote{\CTANurl{ltx-talk}}}{\beamer\footnote{\CTANurl{beamer}}}

\ldeen{Bei Präsentationsfolien spielen (wie bei vielen Texten) eine konzeptionelle Trennung  
  zwischen ihrem Inhalt 
  und  dessen visueller Darstellung. Während sich der Inhalt von Vortrag zu Vortrag ändert, folgt die
  Gestaltung häufig einem einheitlichen Erscheinungsbild oder einem
  vorgegebenen Corporate Design. Dieses Erscheinungsbild wird in der Regel \emph{Thema} genannt.
}{
}\medskip

\ldeen{
  Es ist ausdrücklich \emph{kein} Ziel von \cls{ltx-talk}, Themen zur Verfügung zu stellen.
  Damit bleibt es einzelnen Entwicklern überlassen, Themen für \lttalk{} zu schreiben.
  Bei dieser Arbeit will sie das Paket \lttheme{} ein Stück weit unterstützen.%
  \footnote{Tatsächlich ist \lttheme{} bei dem Versuch entstanden, ein eigenes Thema für
  \lttalk{} im Corporate Design der Universität des Autors zu schreiben.}
  Es stellt ein Theme-System für die Dokumentklasse
  \lttalk{} bereit.  
}{
  It is expressly \emph{not} the goal of \cls{ltx-talk} to provide themes.
  This leaves it up to individual developers to write themes for \lttalk{}.
  The \lttheme{} package is intended to provide some support for this work.%
  \footnote{In fact, \lttheme{} was created while attempting to write a custom theme for
  \lttalk{} in the corporate design of the author's university.}
  It provides a theme engine for the document class \lttalk.
}

\ldeen{
  Ein Thema umfasst dabei mehr als Farben und Schriftarten. Es legt fest,
  welche Informationen dargestellt werden, wie sie auf einer Folie
  angeordnet sind und wie ihre visuelle Gestaltung erfolgt.
  Diese Dinge sind derzeit in \lttalk{} (absichtlich) nur eingeschränkt gestaltbar.
}{
  A theme involves more than just colors and fonts. It determines
  what information is displayed, how it is arranged on a slide,
  and how it is visually designed.
  Currently, these elements can (intentionally) only be customized to a limited extent in \lttalk{}.
}


\begin{warnbox}{\ldeen{Warnung}{Warning}}
  \ldeen{
    Es soll ausdrücklich darauf hingewiesen werden, dass dieses Paket experimentell ist, 
    noch über das Maß hinaus, was häufig experimentelle Pakete haben: Schließlich hängt es von 
    der experimentellen, sich noch in Entwicklung befindenden Klasse \lttalk{} ab. 
    Änderungen dort können zu einer sofortigen Nichtnutzbarkeit von \lttheme{} führen.
  }{
    It should be explicitly noted that this package is experimental, even more so 
    than most experimental packages: After all, it depends on the experimental class \lttalk,
    which is still under development. Changes to that class may immediately render 
    \lttheme{} unusable.
  }
\end{warnbox}

\ldeen{
  Deshalb prüft \lttheme{} beim Laden die unterstützte Version der Klasse.
  Version~0.4.2 unterstützt \cls{ltx-talk}~0.5.3 vom 5.~August~2026.
  Bei einem älteren oder neueren Stand bricht das Paket mit einer gezielten
  Meldung ab. \cls{ltx-talk} bietet gegenwärtig keine eigenen Rollback-Releases
  an; eine bereits mit \cs*{documentclass} geladene Klasse kann das Theme nicht
  nachträglich ersetzen.
}{
  Therefore, \lttheme{} checks the supported class version when it is loaded.
  Version~0.4.2 supports \cls{ltx-talk}~0.5.3 dated 5~August~2026.
  With an older or newer release, the package stops with a targeted diagnostic.
  \cls{ltx-talk} currently provides no rollback releases of its own; a class
  already loaded by \cs*{documentclass} cannot subsequently be replaced by the
  theme.
}

\subsection{\ldeen{Wer sollte diese Dokumentation lesen?}{Who sould read this documentation?}}

\ldeen{
Diese Dokumentation richtet sich an zwei Zielgruppen. Anwender möchten
ein vorhandenes Theme nutzen und gegebenenfalls an ihre Anforderungen
anpassen. Themen-Autoren möchten dagegen ein eigenes Theme erstellen oder
ein vorhandenes Theme erweitern.
}{
This documentation addresses two audiences. Users want to select an
existing theme and, where necessary, adapt it to their requirements.
Theme authors, on the other hand, want to create a new theme or extend
an existing one.
}

\ldeen*{
Für erstere Gruppe ist Teil~@1 gedacht.
}{
  Part~@1 is intended for the former group.
}{\ref{part:user}}
\ldeen*{
Wer dagegen ein eigenes Theme erstellen möchte, sollte anschließend auch
den Teil~@1 lesen. Die Referenz (Teil~@2) kann beiden
Gruppen zum gezielten Nachschlagen dienen.
}{
  On the other hand, anyone who wants to create their own theme should also
  read Part~@1. The reference (Part~@2) can serve as a handy reference for both
  groups.
}{\ref{part:develop-themes}}{\ref{ref:reference}}

\part{\ldeen{Nutzerhandbuch}{User Manual}}
\label{part:user}

\section{\ldeen{Schnelleinstieg}{Quick start}}
\label{sec:quick-start}

\ldeen*{
  Das folgende Beispiel zeigt eine kleine, aber vollständige und 
  lauffähige Präsentation.
  Es lädt das mitgelieferte Theme @1, erzeugt eine Titelfolie
  und eine weitere Folie mit einfachem Inhalt.
}{
  The following example demostrates a mall but complete and executable
  presentation. It loads the provided theme @1, generates a title
  slide and an additional slide with simple content.
}{\code{magpie}}

\begin{example}[
  compile,
  program=lualatex,
  runs=2,
  pages={1,2},
  graphics={width=.46\linewidth},
  pre-output={\begin{mdframed}[
    backgroundcolor=white,
    hidealllines=true,
    innerleftmargin=0pt,
    innerrightmargin=0pt,
    innertopmargin=4pt,
    innerbottommargin=4pt
  ]},
  after-output={\end{mdframed}}
]
  \DocumentMetadata{}
  \documentclass{ltx-talk}

  \usepackage{ltxtalk-theme-magpie}
  \useltxtalktheme{magpie}

  \title{My Presentation}
  \author{Joe Doe}
  \date{\today}

  \begin{document}

  \begin{frame}
    \maketitle
  \end{frame}

  \begin{frame}
    \frametitle{}
    \begin{itemize}
      \item A presentation contains the content.
      \item The theme determines how it is displayed.
    \end{itemize}
  \end{frame}

  \end{document}
\end{example}

\section{\ldeen{Themes verwenden}{Using themes}}
\label{sec:use-themes}
\ldeen*{
Ein Theme wird in zwei Schritten ausgewählt. Zunächst lädt
@1 die Themendatei.
Anschließend aktiviert @2 das Thema für die Präsentation. 
}{
A theme is selected in two steps. First,
@1 loads the theme file.
Then, @2 activates the theme for the presentation. 
}{\cs*{usepackage}\Marg{ltxtalk-theme-magpie}}{\cs*{useltxtalktheme}\Marg{magpie}}

\ldeen{
Die Befehle für Titel, Autor und Datum sowie die Umgebung \code{frame}
gehören zur Dokumentklasse \cls{ltx-talk}. 
Das Thema konfiguriert deren
öffentliche Templates und legt damit fest, wo und in welcher Form diese
Angaben auf den Folien erscheinen. Das Setzen und semantische Taggen der
Inhalte bleibt Aufgabe der Klasse. Der eigentliche Inhalt bleibt dadurch
unabhängig vom gewählten Thema.
}{
% TODO: English translation
}

\ldeen{
Zum Wechseln des Erscheinungsbilds genügt es im einfachsten Fall, ein
anderes mitgeliefertes Theme-Paket zu laden und den zugehörigen Namen an
\cs*{useltxtalktheme} zu übergeben. Welche Themes verfügbar sind und wie
ein Theme auch innerhalb einer Präsentation gewechselt werden kann,
beschreibt der folgende Abschnitt.
}{
% TODO: English translation
}


% - Selecting and switching a theme

\section{\ldeen{Themes anpassen}{Customizing themes}}
\label{sec:customizing-themes}
\subsection{\ldeen{Das Layout anpassen}{Customizing the layout}}

\ldeen{
Nach der Aktivierung lassen sich die vom Theme gewählten Abmessungen mit
\cs*{ltxtalksetup} überschreiben. Dabei werden nur die angegebenen
Eigenschaften geändert. Das folgende Beispiel vergrößert Kopf- und Fußbereich
und schaltet den rechten Randbereich ab:
}{
% TODO: English translation
}

\begin{sourcecode}
  \useltxtalktheme{hawk}
  \ltxtalksetup{
    header-height = 2.8cm,
    footer-height = 0.8cm,
    right-margin = false
  }
\end{sourcecode}

\ldeen{
Zur Verfügung stehen \code{header}, \code{left-margin},
\code{right-margin} und \code{footer} als Schalter sowie
\code{header-height}, \code{footer-height} und \code{margin-width} als
Längen. Eine Änderung der Bereichsgröße verschiebt vorhandene Slots nicht
automatisch. Größere Eingriffe sollten deshalb in einem eigenen Theme oder
Teiltheme vorgenommen werden; Abschnitt~\ref{sec:develop-themes} beschreibt
dieses Vorgehen.
}{
% TODO: English translation
}

\subsection{\ldeen{Farben anpassen}{Customizing colours}}

\ldeen{
Die acht semantischen Theme-Farben lassen sich mit
\cs*{setltxtalkcolors} neu belegen. Jeder Wert besteht wie bei
\cs*{definecolor} aus Farbmodell und Farbangabe:
}{
% TODO: English translation
}

\begin{sourcecode}
  \setltxtalkcolors{
    primary   = {HTML}{005B50},
    secondary = {HTML}{3C7A73},
    accent    = {HTML}{D6A800},
    text      = {HTML}{202020},
    background= {HTML}{FFFFFF},
    structure = {HTML}{005B50},
    alert     = {HTML}{A00020},
    example   = {HTML}{166534}
  }
\end{sourcecode}

\ldeen{
Die entsprechenden Farbnamen tragen das Präfix \code{ltxtalk-}; für
\code{background} lautet der Name abgekürzt \code{ltxtalk-bg}. Die drei
inneren Farben sind zusätzlich unter den von \cls{ltx-talk} verwendeten
Namen \code{structure}, \code{alert} und \code{example} verfügbar. Sie stehen
überall dort zur Verfügung, wo eine Farbe erwartet wird. Ob eine Änderung
sichtbar wird, hängt davon ab, ob das aktive Theme an der betreffenden Stelle
die semantische Farbe oder einen eigenen Farbnamen verwendet.
}{
% TODO: English translation
}

\subsection{\ldeen{Titel platzieren}{Positioning titles}}

\ldeen{
Mit \cs*{setltxtalktitleplacement} wird festgelegt, in welchem Bereich ein
Theme Folien- und Abschnittstitel vorsieht. Ein Folientitel kann im
Inhaltsbereich, im Kopf oder in beiden Bereichen erscheinen; ein
Abschnittstitel kann im Kopf ein- oder ausgeblendet werden:
}{
% TODO: English translation
}

\begin{sourcecode}
  \setltxtalktitleplacement{
    frametitle = both,
    section = header
  }
\end{sourcecode}

\ldeen{
Diese Einstellung ist eine Vereinbarung zwischen Core und Theme. Das Theme
muss passende Slots für die gewählte Platzierung enthalten. Fehlt ein solcher
Slot, erzeugt die Option allein kein zusätzliches sichtbares Element.
}{
% TODO: English translation
}

\subsection{\ldeen{Weitere Anpassungen}{Further customizations}}

\ldeen{
Einzelne vorhandene Slots können nach der Theme-Aktivierung mit
\cs*{styleltxtalkslot} umgestaltet werden. Dafür muss der vom Theme
verwendete Slotname bekannt sein; er gehört deshalb zur Schnittstelle des
jeweiligen Themes, nicht zur allgemeinen Anwenderschnittstelle. Ebenso kann
\cs*{decorateltxtalkarea} beispielsweise eine Linie über dem Footer ändern.
}{
% TODO: English translation
}

\begin{sourcecode}
  \styleltxtalkslot{min-ftitle}{
    color = ltxtalk-primary,
    font = \Large\bfseries,
    align = left
  }
  \decorateltxtalkarea{footer}{
    frame = top,
    rule-color = ltxtalk-accent,
    rule-thickness = 0.6pt
  }
\end{sourcecode}

\ldeen{
Solche Anpassungen sollten unmittelbar nach \cs*{useltxtalktheme} stehen,
weil eine spätere erneute Theme-Aktivierung alle Slots und Dekorationen
zurücksetzt. Für wiederverwendbare oder strukturelle Ergänzungen ist ein
Teiltheme mit \cs*{addltxtalktheme} geeigneter.
}{
% TODO: English translation
}


\section{\ldeen{Barrierefreiheit}{Accessibility}}
\label{sec:accessibility}

\ldeen{
Ein Theme verändert die visuelle Darstellung, darf aber die von
\cls{ltx-talk} erzeugte Dokumentstruktur nicht ersetzen. Folientitel sollten
daher mit \cs*{frametitle}, Untertitel mit \cs*{framesubtitle} und
Gliederungsebenen mit \cs*{section} beziehungsweise \cs*{subsection}
ausgezeichnet werden. Rein optisch formatierter Text kann diese semantischen
Befehle nicht ersetzen.
}{
% TODO: English translation
}

\subsection{\ldeen{Lesbare Gestaltung}{Readable design}}

\ldeen{
Bei eigenen Farbanpassungen müssen Vorder- und Hintergrund ausreichend
kontrastieren. Farbe sollte außerdem nie das einzige Mittel sein, um eine
Bedeutung auszudrücken. Kleine Schrift, sehr dünne Linien und dicht belegte
Randbereiche sind auf Projektoren und bei starker Vergrößerung besonders
problematisch. Entscheidend ist nicht nur die Titelfolie, sondern auch eine
inhaltlich volle Standardfolie.
}{
% TODO: English translation
}

\ldeen{
Text und wesentliche Grafiken sollten innerhalb der verfügbaren Bereiche
bleiben. Eine Änderung von \code{margin-width}, \code{header-height} oder
\code{footer-height} muss deshalb mit langen Titeln, mehrzeiligen Angaben und
verschiedenen Seitenverhältnissen geprüft werden.
}{
% TODO: English translation
}

\subsection{\ldeen{Grafiken und Logos}{Graphics and logos}}

\ldeen{
Für inhaltliche Grafiken gelten die Möglichkeiten von \cls{ltx-talk} und den
verwendeten Grafikpaketen; dort ist ein aussagekräftiger Alternativtext
anzugeben. Wiederkehrende Logos, Farbflächen und Trennlinien sind gewöhnlich
dekorativ und dürfen die Lesereihenfolge nicht mit wiederholten Informationen
belasten. Ein Theme-Autor kennzeichnet solche Elemente durch die passende
Slot-Rolle; Anwender sollten ein inhaltliches Bild daher nicht lediglich als
Theme-Logo einsetzen.
}{
% TODO: English translation
}

\subsection{\ldeen{Tagging und Prüfung}{Tagging and testing}}

\ldeen{
Für getaggte PDF-Ausgabe sollte die Präsentation mit \cs*{DocumentMetadata}
beginnen. Das eigentliche Tagging von Frame-Inhalt und Überschriften übernimmt
gegenwärtig überwiegend \cls{ltx-talk}. \lttheme{} bewahrt diese Anbindung,
indem Folientitel im Inhaltsbereich über die öffentlichen Klassenbefehle
gesetzt werden. Die Accessibility-Angaben der Slot-Schnittstelle dokumentieren
zusätzlich die beabsichtigte Rolle; sie werden noch nicht in jedem Fall in
einen eigenen PDF-Tag des Slots umgesetzt.
}{
% TODO: English translation
}

\ldeen{
Eine Prüfung sollte deshalb sowohl visuell als auch strukturell erfolgen:
mit langen und leeren Metadaten, mehreren Gliederungsebenen, fortgesetzten
Folien und einem PDF-Prüfwerkzeug. Das Bundle enthält dafür l3build-Tests, die
unter anderem die von \cls{ltx-talk} erzeugte Überschriftenstruktur prüfen.
}{
% TODO: English translation
}


% ==========================================================================
% For theme authors
% ==========================================================================


\part{\ldeen{Eigene Themes entwickeln}{Developing custom themes}}
\label{part:develop-themes}

\section{\ldeen{Das Theme-Modell}{The theme model}}
\label{sec:theme-model}

\subsection{\ldeen{Designprinzipien}{Design principles}}
\subsubsection{\ldeen{Semantische Elemente}{Semantic elements}}

\ldeen{
Themes arbeiten mit der Bedeutung von Elementen wie Titel, Autor,
Abschnitt, Logo oder Foliennummer. Sie behandeln diese Elemente nicht
lediglich als Text oder Grafik an einer festen Position.
}{
Themes operate on the meaning of elements such as the title, author,
section, logo, or slide number. They do not treat these elements merely
as text or graphics at fixed positions.
}

\ldeen{
Die technische Repräsentation solcher Elemente durch Slots wird im
Abschnitt über das Theme-Modell erläutert.
}{
The technical representation of such elements by slots is explained
in the section on the theme model.
}


\subsubsection{
  \ldeen
    {Trennung von Anordnung und Erscheinungsbild}
    {Separation of arrangement and appearance}
}

\ldeen{
Die räumliche Anordnung eines Elements und seine visuelle Gestaltung
sind unterschiedliche Eigenschaften. Ein Layout beantwortet vor allem
die Frage, wo und in welcher Größe ein Element erscheint.
}{
The spatial arrangement of an element and its visual styling are
different properties. A layout primarily answers where an element
appears and how much space it occupies.
}

\ldeen{
Farben, Schriften, Linien, Hintergründe und ähnliche Eigenschaften
bestimmen dagegen sein Erscheinungsbild. Durch diese Trennung können
Anordnung und Gestaltung weitgehend unabhängig voneinander entwickelt
werden.
}{
Colours, fonts, rules, backgrounds, and similar properties determine
its appearance. This separation allows arrangement and visual design
to be developed largely independently.
}


\subsubsection{\ldeen{Erweiterbarkeit}{Extensibility}}

\ldeen{
Themes sollen nicht in jedem Fall vollständig neu erstellt werden
müssen. Vorhandene Definitionen können als Ausgangspunkt dienen und
gezielt verändert oder ergänzt werden.
}{
Themes should not always have to be created entirely from scratch.
Existing definitions can serve as a starting point and be changed or
extended selectively.
}

\ldeen{
Auf diese Weise können verwandte Themes gemeinsame Eigenschaften
verwenden und sich nur dort unterscheiden, wo dies tatsächlich
erforderlich ist.
}{
In this way, related themes can share common properties and differ only
where this is actually necessary.
}


\subsubsection{\ldeen{Barrierefreiheit}{Accessibility}}

\ldeen{
Die visuelle Gestaltung einer Präsentation darf ihre semantische
Struktur und ihre Zugänglichkeit nicht unnötig beeinträchtigen.
Barrierefreiheit wird deshalb als Bestandteil des Theme-Modells
behandelt und nicht lediglich nachträglich ergänzt.
}{
The visual design of a presentation should not unnecessarily impair
its semantic structure or accessibility. Accessibility is therefore
treated as part of the theme model rather than being added only as an
afterthought.
}

\ldeen{
Das PDF-Tagging der tatsächlich ausgegebenen Folien übernimmt derzeit
überwiegend die Klasse \cls{ltx-talk}. Ein Theme muss deren semantische
Struktur bewahren; \lttheme{} konfiguriert dafür die öffentlichen Templates
der Klasse. Die im Slot-Modell hinterlegten Angaben zur Barrierefreiheit
beschreiben darüber hinaus die angestrebte Semantik der Slots. Sie sind noch
nicht in allen Fällen unmittelbar mit der von \cls{ltx-talk} erzeugten
Tag-Struktur verbunden.
}{
% TODO: English translation
}

\ldeen{
Die dafür vorgesehenen Eigenschaften und Empfehlungen werden in einem
eigenen Abschnitt beschrieben.
}{
The relevant properties and recommendations are described in a
dedicated section.
}

\subsection{\ldeen{Bereiche}{Areas}}

Ein Theme ordnet Seitenelemente den fünf Bereichen \texttt{header},
\texttt{left}, \texttt{content}, \texttt{right} und \texttt{footer} zu.
Header, Ränder und Footer werden durch den Seitenstil \texttt{lttheme}
ausgegeben. Der Inhaltsbereich bleibt Teil des normalen Frame-Inhalts; dadurch
bleiben Umbruch, Overlays und die semantische Einordnung durch \pkg{ltx-talk}
erhalten.

\subsection{\ldeen{Slots}{Slots}}

Ein Slot bezeichnet ein Element mit einer Aufgabe, nicht eine fest verdrahtete
Stelle. Erst die Zuordnung zu einem Bereich, seine Position und sein
Darstellungs-Template ergeben das konkrete Layout. Mehrere Slots desselben
Typs sind zulässig. So kann etwa ein Folientitel im Header, im Inhaltsbereich
oder in beiden Bereichen vorkommen.

\subsection{\ldeen{Slottypen}{Slot types}}

Vordefinierte Typen liefern unter anderem Dokumenttitel, Autor, Institut,
Abschnitt, Unterabschnitt, Folientitel, Folienuntertitel und Foliennummer.
Der Typ \texttt{custom} bezieht seinen Inhalt aus
\cs{setltxtalkslotcontent}; damit lassen sich beispielsweise Pfad- und
Baumnavigationen oder frei gestaltete Logos ergänzen.

\subsection{\ldeen{Layout}{Layout}}

Die Seitenbereiche werden mit Coffins zusammengesetzt. Jeder Slot besitzt eine
horizontale und vertikale Ankerposition sowie Breite und Höhe. Symbolische
Positionen wie \texttt{left}, \texttt{center}, \texttt{right}, \texttt{top}
und \texttt{bottom} können mit absoluten Längen kombiniert werden. Die
Core-Logik kennt dabei keine Namen konkreter Themes.

\subsection{\ldeen{Gestaltung}{Styling}}

Für jeden Slot existiert eine Instanz des Templates \texttt{lttheme-slot}.
Themes geben geeignete Vorgaben vor; Anwender können diese Instanzen mit den
üblichen Template-Mitteln weiter bearbeiten. Schrift, Farbe, Hintergrund,
Ausrichtung, Innenabstand und Regeln gehören damit zum Theme, nicht zum
Seitenstil.

\subsection{\ldeen{Dekoration von Bereichen}{Area decoration}}

Bereichsdekorationen sind von den Slots unabhängig. Das Template
\texttt{lttheme-area} kann insbesondere Hintergrundflächen und Trennlinien
zeichnen. Eine obere Linie des Footers ist daher eine Eigenschaft des
Footer-Bereichs und kein künstlicher Footer-Slot.

Der Seitenstil bildet die stabile Verbindung zu \lttalk. Sockets teilen
seine Ausgabe zusätzlich in Hintergrund, die vier Seitenbereiche und
Vordergrund. Ein vollständiges Theme muss diese Sockets nicht anfassen; sie
sind Erweiterungspunkte für austauschbare Layoutphasen, nicht der Ersatz für
die additive Slot-Schnittstelle.

\subsection{\ldeen{Aufbau eines Themes}{Structure of a theme}}

\ldeen{
Ein Theme-Paket lädt zunächst \pkg{ltxtalk-theme}, definiert eine
Setup-Anweisung und registriert diese unter einem eindeutigen Namen. Die
Setup-Anweisung wird nicht beim Laden, sondern erst durch
\cs*{useltxtalktheme} ausgeführt. Dadurch können mehrere Theme-Pakete geladen
und später gezielt aktiviert werden.
}{
% TODO: English translation
}

\begin{sourcecode}
  \RequirePackage{ltxtalk-theme}
  \ExplSyntaxOn
  \cs_new_protected:Npn \__my_theme_setup:
    {
      \ltxtalksetup{header=true,header-height=1.5cm}
      % Slots, Positionen und Gestaltung
    }
  \RegisterLTXTalkTheme{my-theme}{\__my_theme_setup:}
  \ExplSyntaxOff
\end{sourcecode}

\ldeen{
Die Core-Logik kennt den Namen des Themes nicht. Bei der Aktivierung leert
sie den bisherigen Theme-Zustand, führt die registrierte Setup-Anweisung aus,
passt den Textbereich an und wählt den Seitenstil \code{lttheme}. Das Theme
beschreibt daher ausschließlich Bereiche, Slots und deren Erscheinungsbild.
}{
% TODO: English translation
}

\subsection{\ldeen{Bereiche und Slots definieren}{Defining areas and slots}}

\ldeen{
Ein Slot wird mit Name, Bereich, Typ und Accessibility-Angaben definiert. Der
Name muss innerhalb des aktiven Themes eindeutig sein. Als Bereiche sind
\code{header}, \code{left}, \code{content}, \code{right} und \code{footer}
vorgesehen.
}{
% TODO: English translation
}

\begin{sourcecode}
  \defineltxtalkslot{my-section}{header}{section}{
    role = navigation
  }
  \defineltxtalkslot{my-path}{right}{custom}{
    role = navigation
  }
  \setltxtalkslotcontent{my-path}{
    \thesection\quad\thesubsection
  }
\end{sourcecode}

\ldeen{
Die vordefinierten Typen \code{title}, \code{author}, \code{date},
\code{institute}, \code{section}, \code{subsection}, \code{frametitle},
\code{framesubtitle} und \code{pagenumber} beziehen ihren Inhalt aus der
Klasse. Der Typ \code{custom} sowie unbekannte eigene Typen verwenden den mit
\cs*{setltxtalkslotcontent} hinterlegten Inhalt. Auch Logos werden zweckmäßig
auf diese Weise mit explizitem Grafikcode belegt.
}{
% TODO: English translation
}

\subsection{\ldeen{Slots positionieren}{Positioning slots}}

\ldeen{
Nach der Definition erhält jeder sichtbare Slot eine Position und eine Größe:
}{
% TODO: English translation
}

\begin{sourcecode}
  \positionltxtalkslot
    {my-section}{left}{center}{0.55\paperwidth}{0.8cm}
  \positionltxtalkslot
    {my-path}{center}{1cm}{2cm}{4cm}
\end{sourcecode}

\ldeen{
Das zweite und dritte Argument bezeichnen die horizontale und vertikale
Position innerhalb des Bereichs. Horizontal sind \code{left}, \code{center}
und \code{right}, vertikal \code{top}, \code{center} und \code{bottom}
zulässig. Stattdessen kann jeweils eine Länge ab der linken beziehungsweise
oberen Bereichskante angegeben werden. Breite und Höhe sind Längen; ein Slot
wird nicht automatisch vor Überlappungen mit anderen Slots geschützt.
}{
% TODO: English translation
}

\subsection{\ldeen{Slots gestalten}{Styling slots}}

\ldeen{
\cs*{styleltxtalkslot} bearbeitet die für den Slot angelegte Instanz des
Templates \code{lttheme-slot}. So bleiben Inhalt, Position und Gestaltung
voneinander getrennt.
}{
% TODO: English translation
}

\begin{sourcecode}
  \styleltxtalkslot{my-section}{
    font = \large\bfseries,
    color = white,
    align = left,
    padding = 2mm,
    frame = bottom,
    rule-color = ltxtalk-accent,
    rule-thick = 0.8pt
  }
\end{sourcecode}

\ldeen{
Neben \code{font}, \code{color}, \code{align} und \code{padding} stehen
\code{background}, \code{frame}, \code{rule-color} und \code{rule-thick}
zur Verfügung. \code{frame} akzeptiert \code{none}, \code{box}, \code{top},
\code{bottom}, \code{left} und \code{right}. Die Schnittstelle kennt außerdem
die Optionen \code{rotate}, \code{opacity}, \code{shadow}, \code{uppercase}
und \code{lowercase}; deren vollständige visuelle Umsetzung ist derzeit noch
nicht abgeschlossen und sollte in stabilen Themes nicht vorausgesetzt werden.
}{
% TODO: English translation
}

\subsection{\ldeen{Bereiche dekorieren}{Decorating areas}}

\ldeen{
Hintergrundflächen und Trennlinien, die einen ganzen Bereich betreffen,
gehören nicht in einen Slot. Sie werden mit \cs*{decorateltxtalkarea} an der
Instanz des Templates \code{lttheme-area} eingestellt:
}{
% TODO: English translation
}

\begin{sourcecode}
  \decorateltxtalkarea{footer}{
    background = white,
    frame = top,
    rule-color = ltxtalk-primary,
    rule-thickness = 0.5pt,
    artifact-type = decorative
  }
\end{sourcecode}

\ldeen{
Für \code{frame} sind zusätzlich \code{topbottom} und die üblichen Seiten
zulässig. Die Dekoration wird vor den Slots des Bereichs gesetzt. Die Optionen
\code{background-opacity}, \code{shadow}, \code{shadow-color},
\code{shadow-offset}, \code{margin-top} und \code{margin-bottom} sind bereits
Teil der Schnittstelle, werden aber noch nicht vollständig gezeichnet.
}{
% TODO: English translation
}

\subsection{
  \ldeen
    {Vorhandene Themes erweitern}
    {Extending existing themes}
}

Ein registriertes Theme kann auch nur zusätzliche Slots oder eine einzelne
Dekoration enthalten. Nach der Aktivierung des Basisthemes fügt
\cs{addltxtalktheme}\marg{Name} ein solches Teiltheme hinzu, ohne vorhandene
Slots zu löschen. Das Beispiel \code{example-switching.tex} ergänzt auf diese
Weise eine eigene Pfadnavigation. Das erneute Aktivieren eines Themes beginnt
dagegen bewusst mit einem leeren Theme-Zustand.

\begin{sourcecode}
  \useltxtalktheme{minimal}
  \addltxtalktheme{course-navigation}
\end{sourcecode}

\ldeen{
Ein Teiltheme sollte nur die Eigenschaften definieren, die es tatsächlich
ergänzt. Insbesondere darf es weder den Theme-Zustand leeren noch selbst den
Seitenstil wechseln. Slotnamen sollten einen paketspezifischen Präfix tragen,
damit sie nicht mit denen des Basisthemes kollidieren.
}{
% TODO: English translation
}

\subsection{\ldeen{Empfehlungen}{Recommendations}}

\ldeen{
Semantische Slottypen sind eigenen Typen vorzuziehen. Dekorative Elemente
sollten als Artefakt, Navigation als solche und echte Überschriften mit einer
passenden Ebene beschrieben werden. Wiederkehrende Seitenelemente dürfen den
eigentlichen Frame-Inhalt und dessen Lesereihenfolge nicht duplizieren.
}{
% TODO: English translation
}

\ldeen{
Themes sollten kurze und lange Inhalte, fehlende Metadaten, mehrzeilige
Kopfzeilen, beide Randbereiche sowie Folien mit und ohne Untertitel abdecken.
Für jedes Theme empfiehlt sich ein l3build-Sichttest; strukturelle Tests sollten
zusätzlich die getaggte Überschriftenhierarchie prüfen. Theme-spezifische
Logik bleibt im Theme-Paket, während der Core nur die allgemeinen Bereiche,
Slots, Templates und den Seitenstil bereitstellt.
}{
% TODO: English translation
}


% ==========================================================================
% Reference
% ==========================================================================

\section{\ldeen{Referenz}{Reference}}
\label{ref:reference}
\subsection{\ldeen{Theme-Verwaltung}{Theme management}}

\begin{commands}
  \command{useltxtalktheme}[\marg{Name}]
  \ldeen{
    Aktiviert das registrierte Theme \meta{Name}. Der bisherige Theme-Zustand
    wird vollständig geleert. \code{default} stellt den Seitenstil der Klasse
    ohne lttheme-Slots wieder her.
  }{% TODO: English translation
  }

  \command{addltxtalktheme}[\marg{Name}]
  \ldeen{
    Führt die Setup-Anweisung eines registrierten Themes zusätzlich zum
    aktiven Theme aus, ohne vorhandene Slots zu löschen. Der Befehl dient vor
    allem zum Hinzufügen von Teilthemes.
  }{% TODO: English translation
  }

  \command{RegisterLTXTalkTheme}[\marg{Name}\marg{Setup-Code}]
  \ldeen{
    Registriert \meta{Setup-Code} unter \meta{Name}. Üblicherweise wird hier
    eine geschützte interne Anweisung des Theme-Pakets übergeben. Eine
    vorhandene Registrierung gleichen Namens wird ersetzt.
  }{% TODO: English translation
  }
\end{commands}

\subsection{\ldeen{Layout-Konfiguration}{Layout configuration}}

\begin{commands}
  \command{ltxtalksetup}[\marg{Optionen}]
  \ldeen{
    Ändert die Geometrie der Theme-Bereiche. Die booleschen Schlüssel
    \code{header}, \code{left-margin}, \code{right-margin} und \code{footer}
    schalten Bereiche ein oder aus. \code{header-height},
    \code{footer-height} und \code{margin-width} erwarten Längen.
    Für asymmetrische Layouts können sichtbare Bereichsbreite und
    Inhaltseinzug getrennt mit \code{left-area-width},
    \code{right-area-width}, \code{content-inset-left} und
    \code{content-inset-right} gesetzt werden.
    \code{header-rows} erwartet eine ganze Zahl; \code{header-columns} und
    \code{footer-columns} erwarten kommagetrennte Anteile. Die beiden letzten
    Spaltenangaben gehören zur älteren Layoutschnittstelle und beeinflussen
    die frei positionierten Slots derzeit nicht.
  }{% TODO: English translation
  }

  \command{DeclareLTXTalkLayout}[\marg{Name}\marg{Optionen}]
  \command{UseLTXTalkLayout}[\marg{Name}]
  \ldeen{
    Registriert eine wiederverwendbare Layoutkonfiguration und aktiviert sie.
    \cs*{UseLTXTalkLayout} darf unmittelbar vor einem \env*{frame} oder als
    erste Anweisung darin stehen. Damit bleiben Theme-Layouts unabhängig von
    der optionalen Argument-Schnittstelle der Klasse.
  }{
    Registers and activates a reusable layout configuration.
    \cs*{UseLTXTalkLayout} may be used immediately before a \env*{frame} or
    as its first command. Theme layouts therefore remain independent of the
    class's optional frame argument interface.
  }

  \command{LTXTalkBaseLeftMargin}
  \command{LTXTalkBaseRightMargin}
  \command{LTXTalkCurrentLeftMargin}
  \command{LTXTalkCurrentRightMargin}
  \ldeen{
    Expandierbare Längen für die ursprünglichen beziehungsweise aktuell
    wirksamen Seitenränder. Theme-Code kann damit ohne Zugriff auf interne
    Variablen an der Textkante ausrichten.
  }{
    Expandable lengths for the original and currently effective page margins.
    Theme code can use them to align content without accessing internals.
  }
\end{commands}

\subsection{\ldeen{Metadaten}{Metadata}}

\begin{commands}
  \command{LTXTalkMetadata}[\marg{Feld}]
  \command{UseLTXTalkMetadata}[\oarg{Optionen}\marg{Feld}]
  \ldeen{
    Liefert Präsentationsmetadaten über eine öffentliche Schnittstelle.
    Verfügbare Felder sind \code{title}, \code{short-title}, \code{subtitle},
    \code{short-subtitle}, \code{author}, \code{short-author},
    \code{institute}, \code{short-institute}, \code{date},
    \code{short-date}, \code{section}, \code{subsection},
    \code{frametitle}, \code{framesubtitle}, \code{framenumber} und
    \code{totalframes}. Die zweite Form akzeptiert \code{lines=1}, um
    erzwungene Zeilenumbrüche durch \code{linebreak} (standardmäßig ein
    Leerzeichen) zu ersetzen.
  }{
    Provides presentation metadata through a public interface. Available
    fields are \code{title}, \code{short-title}, \code{subtitle},
    \code{short-subtitle}, \code{author}, \code{short-author},
    \code{institute}, \code{short-institute}, \code{date},
    \code{short-date}, \code{section}, \code{subsection},
    \code{frametitle}, \code{framesubtitle}, \code{framenumber}, and
    \code{totalframes}. The second form accepts \code{lines=1} to replace
    forced line breaks with \code{linebreak}, which defaults to a space.
  }
\end{commands}

\subsection{\ldeen{Farbkonfiguration}{Colour configuration}}

\begin{commands}
  \command{setltxtalkcolors}[\marg{Farben}]
  \ldeen{
    Definiert die semantischen Farben \code{primary}, \code{secondary},
    \code{accent}, \code{text}, \code{background}, \code{structure},
    \code{alert} und \code{example}. Jeder Wert hat die Form
    \code{\Marg{Modell}\Marg{Wert}}, beispielsweise
    \code{primary=\Marg{HTML}\Marg{005B50}}. Die zugehörigen Farbnamen lauten
    \code{ltxtalk-primary}, \code{ltxtalk-secondary}, \code{ltxtalk-accent},
    \code{ltxtalk-text}, \code{ltxtalk-bg}, \code{ltxtalk-structure},
    \code{ltxtalk-alert} und \code{ltxtalk-example}. Die letzten drei werden
    zugleich auf die Klassennamen \code{structure}, \code{alert} und
    \code{example} abgebildet.
  }{% TODO: English translation
  }
\end{commands}

\subsection{\ldeen{Inneres Theme}{Inner theme}}

\begin{commands}
  \command{styleltxtalkblock}[\marg{Typ}\marg{Optionen}]
  \ldeen{
    Bearbeitet die Instanz \code{block}, \code{alertblock} oder
    \code{exampleblock} des Template-Typs \code{lttheme-block}. Die
    Standard-Schnittstelle umfasst \code{background}, \code{rule-color},
    \code{rule-thickness}, \code{padding}, \code{corner-radius},
    \code{shadow}, \code{shadow-offset}, \code{before} und \code{after}.
    Eigene Template-Codes können weitere Darstellungen, etwa echte Rundungen
    oder Schatten, ergänzen. Die Standardinstanz bleibt ohne Rahmen und
    bewahrt damit die Darstellung der Klasse.
  }{% TODO: English translation
  }

  \command{setltxtalkitemize}[\marg{Optionen}]
  \ldeen{
    Setzt mit \code{level-1} bis \code{level-4} das Aufzählungszeichen der
    jeweiligen \code{itemize}-Ebene. Es wird nur die visuelle Marke geändert;
    Listen- und Listenelement-Tags der Klasse bleiben unverändert.
  }{% TODO: English translation
  }
\end{commands}

\subsection{\ldeen{Slot-Definition}{Slot definition}}

\begin{commands}
  \command{defineltxtalkslot}
    [\marg{Name}\marg{Bereich}\marg{Typ}\marg{Accessibility-Optionen}]
  \ldeen{
    Definiert einen Slot und fügt ihn am Ende des angegebenen Bereichs ein.
    Gültige Standardbereiche sind \code{header}, \code{left}, \code{content},
    \code{right} und \code{footer}. Vordefinierte Typen sind \code{title},
    \code{author}, \code{date}, \code{institute}, \code{section},
    \code{subsection}, \code{frametitle}, \code{framesubtitle},
    \code{pagenumber}, \code{logo} und \code{custom}.
  }{% TODO: English translation
  }

  \command{setltxtalkslotcontent}[\marg{Name}\marg{Inhalt}]
  \ldeen{
    Hinterlegt den auszuführenden Inhalt eines freien Slots. Er wird für
    \code{custom}, \code{logo} und nicht vordefinierte eigene Typen verwendet.
  }{% TODO: English translation
  }
\end{commands}

\subsection{\ldeen{Slot-Positionierung}{Slot positioning}}

\begin{commands}
  \command{positionltxtalkslot}
    [\marg{Name}\marg{x}\marg{y}\marg{Breite}\marg{Höhe}]
  \ldeen{
    Positioniert den Slot relativ zu seinem Bereich. Für \meta{x} sind
    \code{left}, \code{center}, \code{right} oder eine Länge zulässig; für
    \meta{y} entsprechend \code{top}, \code{center}, \code{bottom} oder eine
    Länge. Positive Längen werden von der linken beziehungsweise oberen Kante
    gemessen.
  }{% TODO: English translation
  }

  \command{DeclareLTXTalkAreaGrid}
    [\marg{Bereich}\marg{Spalten}\marg{Zeilen}]
  \command{PositionLTXTalkSlotInGrid}
    [\marg{Name}\marg{Spalte}\marg{Zeile}\marg{Spaltenbreite}\marg{Zeilenhöhe}]
  \command{SetLTXTalkAreaGridRows}[\marg{Bereich}\marg{Zeilen}]
  \ldeen{
    Definiert ein Raster für einen Bereich und positioniert Slots darin.
    Spalten sind kommagetrennte Längen; \code{*} verteilt den verbleibenden
    Platz gleichmäßig auf alle Sternspalten. Spalten und Zeilen werden ab eins
    gezählt. Der letzte Befehl ändert die Zeilenzahl lokal und eignet sich
    daher für Umschalter am Anfang eines Frames.
  }{
    Defines a grid for an area and positions slots in it. Columns are a
    comma-separated list of lengths; \code{*} distributes the remaining
    width equally among all star columns. Columns and rows are one-based. The
    final command changes the row count locally and is therefore suitable for
    switches at the beginning of a frame.
  }
\end{commands}

\subsection{\ldeen{Titelseiten}{Title pages}}

\begin{commands}
  \command{DeclareLTXTalkTitleLayout}
    [\marg{Name}\marg{Layoutoptionen}\marg{Inhalt}]
  \command{UseLTXTalkTitleLayout}[\marg{Name}]
  \command{MakeLTXTalkTitle}[\oarg{Name}]
  \ldeen{
    Registriert eine Titelseitenvariante aus temporären Layoutoptionen und
    ihrem Inhalt. \cs*{UseLTXTalkTitleLayout} wählt die Vorgabe;
    \cs*{MakeLTXTalkTitle} erzeugt den Frame und kann optional eine Variante
    nur für diesen Aufruf wählen. Nach der Titelseite wird die vorherige
    Geometrie wiederhergestellt.
  }{
    Registers a title-page variant consisting of temporary layout options and
    its contents. \cs*{UseLTXTalkTitleLayout} selects the default;
    \cs*{MakeLTXTalkTitle} creates the frame and can optionally select a
    variant for that invocation. The previous geometry is restored after the
    title page.
  }
  \command{LTXTalkHeaderRows}
  \command{IfLTXTalkTitlePageTF}[\marg{Titelseitencode}\marg{Framecode}]
  \ldeen{
    Liefert Theme-Autoren die aktuell wirksame Zahl der Kopfzeilen bzw.
    unterscheidet eine von \cs*{MakeLTXTalkTitle} erzeugte Titelseite von
    einem gewöhnlichen Frame. Beide Schnittstellen bleiben auch dann gültig,
    wenn die Basisklasse den Seitenkopf erst beim Ausschiffen aufbaut.
  }{
    Gives theme authors the effective number of header rows and distinguishes
    a title page created by \cs*{MakeLTXTalkTitle} from an ordinary frame.
    Both interfaces remain valid when the base class defers construction of
    the running head until shipout.
  }
\end{commands}

\subsection{\ldeen{Slot-Gestaltung}{Slot styling}}

\begin{commands}
  \command{styleltxtalkslot}[\marg{Name}\marg{Optionen}]
  \ldeen{
    Bearbeitet die Template-Instanz eines Slots. Stabil umgesetzt sind
    \code{font}, \code{color}, \code{background}, \code{align}, \code{frame},
    \code{padding}, \code{rule-color} und \code{rule-thick}. \code{align}
    akzeptiert \code{left}, \code{center} und \code{right}; \code{frame}
    akzeptiert \code{none}, \code{box}, \code{top}, \code{bottom},
    \code{left} und \code{right}.
  }{% TODO: English translation
  }
\end{commands}

\subsection{\ldeen{Bereichsdekoration}{Area decoration}}

\begin{commands}
  \command{decorateltxtalkarea}[\marg{Bereich}\marg{Optionen}]
  \ldeen{
    Bearbeitet die Template-Instanz eines Bereichs. Stabil umgesetzt sind
    \code{background}, \code{frame}, \code{rule-color},
    \code{rule-thickness}, \code{padding} und \code{artifact-type}.
    \code{frame} akzeptiert \code{none}, \code{box}, \code{top},
    \code{bottom}, \code{left}, \code{right} und \code{topbottom}.
  }{% TODO: English translation
  }
\end{commands}

\subsection{\ldeen{Titelplatzierung}{Title placement}}

\begin{commands}
  \command{setltxtalktitleplacement}[\marg{Optionen}]
  \ldeen{
    Legt mit \code{frametitle=content|header|both} und
    \code{section=header|none} die beabsichtigte Titelplatzierung fest. Das
    aktive Theme muss Slots für die gewählten Bereiche definieren.
  }{% TODO: English translation
  }
\end{commands}

\subsection{
  \ldeen
    {Optionen für Barrierefreiheit}
    {Accessibility options}
}

\ldeen{
Das vierte Argument von \cs*{defineltxtalkslot} akzeptiert folgende
Schlüssel: \code{role} mit den Werten \code{heading}, \code{heading-1},
\code{heading-2}, \code{heading-3}, \code{artifact}, \code{content},
\code{navigation} oder \code{note}; \code{heading-level} mit einer ganzen
Zahl; \code{alt-text}; \code{continuation-text}; sowie
\code{continuation-mode=explicit|implicit}. Diese Angaben beschreiben die
beabsichtigte Semantik des Slots. Ihre Abbildung auf PDF-Tags ist derzeit noch
nicht für alle Slots implementiert.
}{
% TODO: English translation
}

\begin{commands}
  \command{setltxtalkheadingbehavior}[\marg{Optionen}]
  \ldeen{
    Konfiguriert die Überschriftenbehandlung mit den booleschen Schlüsseln
    \code{auto-section-title}, \code{detect-continuation} und
    \code{header-as-navigation}. Diese Schnittstelle ist experimentell; das
    tatsächlich erzeugte Tagging wird wesentlich durch \cls{ltx-talk}
    bestimmt.
  }{% TODO: English translation
  }
\end{commands}

\subsection{\ldeen{Templates und Sockets}{Templates and sockets}}

\ldeen{
Für jeden Slot existiert eine Instanz des Template-Typs
\code{lttheme-slot}, für jeden Bereich eine Instanz von \code{lttheme-area}.
Zum Bearbeiten dieser Instanzen dienen die beiden zuvor dokumentierten
Befehle für Slotgestaltung und Bereichsdekoration.
}{
% TODO: English translation
}

\ldeen{
Der Seitenstil stellt Sockets für den Seitenhintergrund, die Bereiche
\code{header}, \code{left}, \code{right} und \code{footer} sowie den
Seitenvordergrund bereit. Ihre Namen folgen dem Schema
\code{lttheme/area/}\meta{Bereich} beziehungsweise
\code{lttheme/page/}\meta{Ebene}. Der Plug \code{standard} rendert den
zugehörigen Bereich. Eigene Plugs sind für grundlegende, austauschbare
Ausgabephasen gedacht; zusätzliche Seitenelemente sollten gewöhnlich als
Slots eines Teilthemes implementiert werden.
}{
% TODO: English translation
}

\begin{table}[ht]
\centering
\begin{tabular}{lll}
\hline
\textbf{Type} & \textbf{Description} & \textbf{Default Role} \\
\hline
\texttt{frametitle}    & Frame title        & heading (H2)     \\
\texttt{framesubtitle} & Frame subtitle     & heading (H3)     \\
\texttt{title}         & Presentation title & heading (H1)     \\
\texttt{section}       & Current section    & heading (H2)     \\
\texttt{subsection}    & Current subsection & heading (H3)     \\
\texttt{author}        & Author name        & note             \\
\texttt{date}          & Date               & note             \\
\texttt{institute}     & Institute name     & note             \\
\texttt{pagenumber}    & Page number        & navigation       \\
\texttt{logo}          & Logo               & graphic artifact \\
\texttt{custom}        & User-defined       & artifact         \\
\hline
\end{tabular}
\caption{Available Template Variables}
\label{tab:template-variables}
\end{table}

\begin{table}[ht]
\centering
\begin{tabular}{ll}
\hline
\textbf{Area} & \textbf{Description} \\
\hline
\texttt{header}  & Top of the frame     \\
\texttt{left}    & Left margin          \\
\texttt{content} & Main content area    \\
\texttt{right}   & Right margin         \\
\texttt{footer}  & Bottom of the frame  \\
\hline
\end{tabular}
\caption{Available Areas}
\label{tab:available-areas}
\end{table}

\part{\ldeen{Referenz}{Reference}}
\label{part:reference}

% bisherige API-Referenz


%\section{\ldeen{Nutzung}{Usage}}
%
% \section{Overview}
% The \textsf{lttheme} module provides a theme engine and four predefined
% themes for the \textsf{ltx-talk} document class.
%
% \section{API reference}
% The following material was migrated from \texttt{docs/api-reference.md}.
%
% # ltxtalk-themes API Reference
% 
% ## Overview
% 
% The ltxtalk-themes system provides a flexible, accessible theme engine for the ltx-talk document class. Themes are built on a slot-based architecture using LaTeX3 coffins internally, but this complexity is completely hidden from users.
% 
% ## Core Commands
% 
% ### Theme Management
% 
% | Command | Description |
% |---------|-------------|
% | `\useltxtalktheme{name}` | Activate a theme |
% | `\RegisterLTXTalkTheme{name}{setup}` | Register a new theme (for developers) |
% 
% ### Layout Configuration
% 
% ```latex
% \ltxtalksetup{
%   header = true/false,
%   header-height = dimension,
%   header-rows = integer,
%   left-margin = true/false,
%   right-margin = true/false,
%   footer = true/false,
%   footer-height = dimension,
%   margin-width = dimension,
% }
% ```
% 
% ### Color Configuration
% 
% ```latex
% \setltxtalkcolors{
%   primary = {HTML}{2B3A67},
%   secondary = {HTML}{496A81},
%   accent = {HTML}{66999B},
%   text = {HTML}{333333},
%   background = {HTML}{FFFFFF},
% }
% ```
% 
% ### Slot Definition
% ```latex
% \defineltxtalkslot{slot-name}{area}{type}{accessibility-options}
% 
% % Example:
% \defineltxtalkslot{my-title}{header}{frametitle}{
%   role = heading,
%   heading-level = 2
% }
% ```
% 
% ### Slot Positioning
% ```latex
% \positionltxtalkslot{slot-name}{x-pos}{y-pos}{width}{height}
% 
% % Example:
% \positionltxtalkslot{my-title}{center}{0.5cm}{0.8\paperwidth}{1.5cm}
% ```
% 
% ### Slot Styling
% ```latex
% \styleltxtalkslot{slot-name}{
%   font = \Large\bfseries,
%   color = ltxtalk-primary,
%   frame = bottom,
%   rule-color = ltxtalk-accent,
%   rule-thick = 2pt,
%   padding = 5pt,
%   align = center,
%   background = none,
% }
% ```
% 
% ### Area Decoration
% ```latex
% \decorateltxtalkarea{header}{
%   frame = bottom,
%   rule-color = ltxtalk-accent,
%   rule-thickness = 4pt,
%   background = ltxtalk-primary!5,
%   shadow = true,
% }
% ```
% 
% ### Available Slot Types
% 
% | Type            | Description        | Default Role.    |
% |-----------------|--------------------|------------------| 
% | `frametitle`    | Frame title        | heading (H2)     | 
% | `framesubtitle` | Frame subtitle     | heading (H3)     | 
% | `title`         | Presentation title | heading (H1)     | 
% | `section`       | Current section    | heading (H2)     | 
% | `subsection`    | Current subsection | heading (H3)     | 
% | `author`        | Author name        | note             | 
% | `date`          | Date               | note             | 
% | `institute`     | Institute name     | note             | 
% | `pagenumber`    | Page number        | navigation       | 
% | `logo`          | Logo               | graphic artifact |
% | `custom`        | User-defined       | artifact         | 
% 
% ### Available Areas
% 
% | Area.     | Description           |
% |-----------|-----------------------| 
% | `header`  | Top of the frame      | 
% | `left`.   | Left margin           | 
% | `content` | Main content area     | 
% | `right`.  | Right margin          | 
% | `footer`  | Bottom of the frame.  | 
% 
% ### Accessibility Options
% ```latex
% \defineltxtalkslot{name}{area}{type}{
%   role = heading|artifact|navigation|content|note,
%   heading-level = 1|2|3,
%   alt-text = "Description",
%   continuation-text = "(cont.)",
%   continuation-mode = explicit|implicit,
% }
% ```
% 
% ### Title Placement
% ```latex
% \setltxtalktitleplacement{
%   frametitle = content|header|both,
%   section = header|none,
% }
% ```
% 
% ### Best Practices
% 1. Always provide accessibility roles for slots
% 2. Use `artifact` role for decorative elements
% 3. Set appropriate heading levels for titles
% 4. Test with screen readers
% 5. Provide alt-text for logos and images
% \end{lang}
%
% \begin{lang}{de}
% \section{Überblick}
% Das Modul \textsf{lttheme} stellt eine Theme-Engine und vier
% vordefinierte Themes für die Dokumentklasse \textsf{ltx-talk} bereit.
% \section{API-Referenz}
% Die englische API-Referenz wurde zunächst unverändert übernommen.
% \end{lang}
%
\part{Implementation}
\section{lttheme.sty}
%\AutoDocInput{lttheme.dtx}{pkg}
\section{\ldeen{Enthaltene Themen}{Includes Themes}}
%\AutoDocInput{lttheme.dtx}{minimal}
%\AutoDocInput{lttheme.dtx}{magpie}
%\AutoDocInput{lttheme.dtx}{hawk}
%\AutoDocInput{lttheme.dtx}{sparrow}
%\section{\ldeen{Farbthemen}{Color Themes}}
%\AutoDocInput{lttheme.dtx}{anchovy}
%\AutoDocInput{lttheme.dtx}{carp}
\end{document}
%</driver>
%<*pkg>
\NeedsTeXFormat{LaTeX2e}
\ProvidesExplPackage{ltxtalk-theme}
{2026/08/10}
{v0.4.2}
{Theme engine for ltx-talk}

\RequirePackage{xcoffins}
\RequirePackage{xcolor}

\msg_new:nnnn { ltxtalk } { missing-class }
  { Missing~document~class }
  { The~ltxtalk-theme~package~requires~ltx-talk. }
\msg_new:nnnn { ltxtalk } { class-too-old }
  { Unsupported~ltx-talk~version }
  {
    ltxtalk-theme~v0.4.2~requires~ltx-talk~v0.5.3~from~2026-08-05.
    Update~the~ltx-talk~document~class~before~loading~this~theme~engine.
  }
\msg_new:nnnn { ltxtalk } { class-too-new }
  { Untested~ltx-talk~version }
  {
    ltxtalk-theme~v0.4.2~was~tested~with~ltx-talk~v0.5.3~from~2026-08-05.
    The~loaded~class~is~newer.~Install~a~compatible~theme-engine~release~or~
    update~its~declared~compatibility~range~after~running~the~test~suite.
  }

% \ldeen*{Das Paket ist nur sinnvoll, wenn @1 geladen ist.}{
%  The package is only useful, if @1 is loaded. 
% }{\cls*{ltx-talk}}
\@ifclassloaded{ltx-talk}
  {
    \@ifclasslater{ltx-talk}{2026/08/05}
      {}
      {\msg_fatal:nn { ltxtalk } { class-too-old }}
    \@ifclasslater{ltx-talk}{2026/08/06}
      {\msg_fatal:nn { ltxtalk } { class-too-new }}
      {}
  }
  { 
    \msg_error:nnnn { ltxtalk } { missing-class }
      { This~package~requires~the~ltx-talk~document~class }
      { Please~use~\token_to_str:N \documentclass{ltx-talk} } 
  }

% ==========================================
% NACHRICHTEN
% ==========================================
\msg_new:nnnn { ltxtalk } { unknown-theme }
  { Unknown~theme~'#1'. }
  { The~theme~'#1'~is~not~registered.~Use~\token_to_str:N \RegisterLTXTalkTheme~to~register~it. }

\msg_new:nnn { ltxtalk } { missing-heading }
  { Warning:~No~heading~defined~in~frame.~This~may~affect~accessibility. }

\msg_new:nnn { ltxtalk } { theme-not-activated }
  { Theme~'#1'~loaded~but~not~activated.~Use~\token_to_str:N \useltxtalktheme{#1}~to~activate. }
\msg_new:nnnn { ltxtalk } { unknown-layout }
  { Unknown~layout~'#1'. }
  { Declare~it~with~\token_to_str:N\DeclareLTXTalkLayout. }
\msg_new:nnnn { ltxtalk } { unknown-title-layout }
  { Unknown~title~layout~'#1'. }
  { Declare~it~with~\token_to_str:N\DeclareLTXTalkTitleLayout. }

% ==========================================
% INTERNE COFFINS (VERSTECKT VOR NUTZER)
% ==========================================
\coffin_new:N \g__ltxtalk_page_coffin
\coffin_new:N \g__ltxtalk_header_coffin
\coffin_new:N \g__ltxtalk_left_coffin
\coffin_new:N \g__ltxtalk_content_coffin
\coffin_new:N \g__ltxtalk_right_coffin
\coffin_new:N \g__ltxtalk_footer_coffin

% Temporäre Coffins für Spalten-Layout
\coffin_new:N \l__ltxtalk_header_left_col
\coffin_new:N \l__ltxtalk_header_center_col
\coffin_new:N \l__ltxtalk_header_right_col

% ==========================================
% THEME-REGISTRIERUNG
% ==========================================

\tl_new:N \g__ltxtalk_active_theme_tl
\tl_gset:Nn \g__ltxtalk_active_theme_tl { default }

\prop_new:N \g__ltxtalk_themes_prop

\NewDocumentCommand \RegisterLTXTalkTheme { m m }
  {
    \prop_gput:Nnn \g__ltxtalk_themes_prop { #1 } { #2 }
  }

\NewDocumentCommand \useltxtalktheme { m }
  {
    \str_case:nnF { #1 }
      {
        { default } { \__ltxtalk_activate_default_theme: }
      }
      {
        \prop_get:NnNTF \g__ltxtalk_themes_prop { #1 } \l_tmpa_tl
          {
            \tl_gset:Nn \g__ltxtalk_active_theme_tl { #1 }
            \l_tmpa_tl
            \__ltxtalk_apply_class_templates:
          }
          {
            \msg_error:nnn { ltxtalk } { unknown-theme } { #1 }
          }
      }
  }

\cs_new_protected:Npn \__ltxtalk_activate_default_theme:
  {
    \tl_gset:Nn \g__ltxtalk_active_theme_tl { default }
  }

% ==========================================
% LAYOUT-KONFIGURATION
% ==========================================

\dim_new:N \l__ltxtalk_margin_dim
\dim_new:N \l__ltxtalk_left_area_width_dim
\dim_new:N \l__ltxtalk_right_area_width_dim
\dim_new:N \l__ltxtalk_content_inset_left_dim
\dim_new:N \l__ltxtalk_content_inset_right_dim
\dim_new:N \l__ltxtalk_header_height_dim
\dim_new:N \l__ltxtalk_footer_height_dim
\dim_new:N \l__ltxtalk_area_width_dim
\dim_new:N \l__ltxtalk_area_height_dim

\int_new:N \l__ltxtalk_header_rows_int
\int_new:N \g__ltxtalk_heading_id_int

\bool_new:N \l__ltxtalk_show_header_bool
\bool_new:N \l__ltxtalk_show_left_margin_bool
\bool_new:N \l__ltxtalk_show_right_margin_bool
\bool_new:N \l__ltxtalk_show_footer_bool
\bool_new:N \l__ltxtalk_frametitle_in_header_bool
\bool_new:N \l__ltxtalk_explicit_cont_bool
\bool_new:N \l__ltxtalk_auto_section_title_bool
\bool_new:N \l__ltxtalk_detect_cont_bool
\bool_new:N \l__ltxtalk_header_nav_bool

\clist_new:N \l__ltxtalk_header_columns_clist
\clist_new:N \l__ltxtalk_footer_columns_clist
\prop_new:N \g__ltxtalk_layouts_prop

\keys_define:nn { ltxtalk / layout }
  {
    header         .bool_set:N = \l__ltxtalk_show_header_bool,
    header         .initial:n = false,
    header-height  .dim_set:N = \l__ltxtalk_header_height_dim,
    header-height  .initial:n = 1.5cm,
    header-rows    .int_set:N = \l__ltxtalk_header_rows_int,
    header-rows    .initial:n = 1,
    left-margin    .bool_set:N = \l__ltxtalk_show_left_margin_bool,
    left-margin    .initial:n = false,
    right-margin   .bool_set:N = \l__ltxtalk_show_right_margin_bool,
    right-margin   .initial:n = false,
    footer         .bool_set:N = \l__ltxtalk_show_footer_bool,
    footer         .initial:n = false,
    footer-height  .dim_set:N = \l__ltxtalk_footer_height_dim,
    footer-height  .initial:n = 1cm,
    margin-width   .code:n =
      {
        \dim_set:Nn \l__ltxtalk_margin_dim {#1}
        \dim_set:Nn \l__ltxtalk_left_area_width_dim {#1}
        \dim_set:Nn \l__ltxtalk_right_area_width_dim {#1}
        \dim_set:Nn \l__ltxtalk_content_inset_left_dim {#1}
        \dim_set:Nn \l__ltxtalk_content_inset_right_dim {#1}
      },
    margin-width   .initial:n = 2.5cm,
    left-area-width .dim_set:N = \l__ltxtalk_left_area_width_dim,
    left-area-width .initial:n = 2.5cm,
    right-area-width .dim_set:N = \l__ltxtalk_right_area_width_dim,
    right-area-width .initial:n = 2.5cm,
    content-inset-left .dim_set:N = \l__ltxtalk_content_inset_left_dim,
    content-inset-left .initial:n = 2.5cm,
    content-inset-right .dim_set:N = \l__ltxtalk_content_inset_right_dim,
    content-inset-right .initial:n = 2.5cm,
    
    header-columns .clist_set:N = \l__ltxtalk_header_columns_clist,
    header-columns .initial:n = { 0.2, 0.6, 0.2 },
    footer-columns .clist_set:N = \l__ltxtalk_footer_columns_clist,
    footer-columns .initial:n = { 0.33, 0.34, 0.33 },
  }

\NewDocumentCommand \ltxtalksetup { m }
  {
    \keys_set:nn { ltxtalk / layout } { #1 }
    \__ltxtalk_update_text_area:
  }

\NewDocumentCommand \DeclareLTXTalkLayout { m m }
  { \prop_gput:Nnn \g__ltxtalk_layouts_prop {#1} {#2} }

\NewDocumentCommand \UseLTXTalkLayout { m }
  {
    \prop_get:NnNTF \g__ltxtalk_layouts_prop {#1} \l_tmpa_tl
      { \keys_set:nV {ltxtalk/layout} \l_tmpa_tl \__ltxtalk_update_text_area: }
      { \msg_error:nnn {ltxtalk}{unknown-layout}{#1} }
  }

\prop_new:N \g__ltxtalk_title_layout_options_prop
\prop_new:N \g__ltxtalk_title_layout_code_prop
\tl_new:N \g__ltxtalk_title_layout_tl
\bool_new:N \l__ltxtalk_title_page_bool
\int_new:N \g__ltxtalk_title_header_rows_int
\int_new:N \g__ltxtalk_saved_header_grid_rows_int

\NewExpandableDocumentCommand \LTXTalkHeaderRows {}
  {\bool_if:NTF \l__ltxtalk_title_page_bool
     {\int_use:N \g__ltxtalk_title_header_rows_int}
     {\int_use:N \l__ltxtalk_header_rows_int}}

\NewDocumentCommand \IfLTXTalkTitlePageTF {m m}
  {\bool_if:NTF \l__ltxtalk_title_page_bool {#1}{#2}}

\NewDocumentCommand \DeclareLTXTalkTitleLayout { m m m }
  {
    \prop_gput:Nnn \g__ltxtalk_title_layout_options_prop {#1} {#2}
    \prop_gput:Nnn \g__ltxtalk_title_layout_code_prop {#1} {#3}
  }

\NewDocumentCommand \UseLTXTalkTitleLayout { m }
  { \tl_gset:Nn \g__ltxtalk_title_layout_tl {#1} }

\NewDocumentCommand \MakeLTXTalkTitle {O{}}
  {
    \tl_if_blank:nTF {#1}
      {\tl_set_eq:NN \l_tmpa_tl \g__ltxtalk_title_layout_tl}
      {\tl_set:Nn \l_tmpa_tl {#1}}
    \prop_get:NVNTF \g__ltxtalk_title_layout_code_prop \l_tmpa_tl \l_tmpb_tl
      {
        \prop_get:NVN \g__ltxtalk_title_layout_options_prop \l_tmpa_tl \l_tmpc_tl
        \group_begin:
          \keys_set:nV {ltxtalk/layout} \l_tmpc_tl
          % ltx-talk 0.5.3 can defer construction of the running head until
          % after the frame group.  Keep the small amount of state needed by
          % header sockets alive until this page has actually shipped.
          \bool_gset_true:N \l__ltxtalk_title_page_bool
          \int_gset_eq:NN \g__ltxtalk_title_header_rows_int
            \l__ltxtalk_header_rows_int
          \prop_if_exist:NT \g__ltxtalk_area_grid_rows_prop
            {
              \prop_get:NnN \g__ltxtalk_area_grid_rows_prop {header} \l_tmpa_tl
              \int_gset:Nn \g__ltxtalk_saved_header_grid_rows_int {\l_tmpa_tl}
              \prop_gput:NnV \g__ltxtalk_area_grid_rows_prop {header}
                \l__ltxtalk_header_rows_int
            }
          \AddToHookNext{shipout/after}
            {
              \bool_gset_false:N \l__ltxtalk_title_page_bool
              \prop_if_exist:NT \g__ltxtalk_area_grid_rows_prop
                {\prop_gput:Nnx \g__ltxtalk_area_grid_rows_prop {header}
                  {\int_use:N \g__ltxtalk_saved_header_grid_rows_int}}
            }
          \__ltxtalk_update_text_area:
          \begin{frame}[vertical-alignment=top]
            \thispagestyle{lttheme}
            \tl_use:N \l_tmpb_tl
          \end{frame}
        \group_end:
        \__ltxtalk_update_text_area:
      }
      {\msg_error:nnV {ltxtalk}{unknown-title-layout}\l_tmpa_tl}
  }

% ==========================================
% FARBSCHEMA
% ==========================================

\keys_define:nn { ltxtalk / colors }
  {
    primary    .code:n = \__ltxtalk_set_color:nn { ltxtalk-primary } { #1 },
    secondary  .code:n = \__ltxtalk_set_color:nn { ltxtalk-secondary } { #1 },
    accent     .code:n = \__ltxtalk_set_color:nn { ltxtalk-accent } { #1 },
    text       .code:n = \__ltxtalk_set_color:nn { ltxtalk-text } { #1 },
    background .code:n = \__ltxtalk_set_color:nn { ltxtalk-bg } { #1 },
    structure  .code:n = \__ltxtalk_set_semantic_color:nn { structure } { #1 },
    alert      .code:n = \__ltxtalk_set_semantic_color:nn { alert } { #1 },
    example    .code:n = \__ltxtalk_set_semantic_color:nn { example } { #1 },
  }

\cs_new_protected:Npn \__ltxtalk_set_color:nn #1#2
  { \__ltxtalk_set_color_aux:nnn { #1 } #2 }

\cs_new_protected:Npn \__ltxtalk_set_color_aux:nnn #1#2#3
  { \color_set:nnn { #1 } { #2 } { #3 } }

\cs_new_protected:Npn \__ltxtalk_set_semantic_color:nn #1#2
  { \__ltxtalk_set_semantic_color_aux:nnn { #1 } #2 }
\cs_new_protected:Npn \__ltxtalk_set_semantic_color_aux:nnn #1#2#3
  {
    \color_set:nnn { ltxtalk-#1 } { #2 } { #3 }
    \color_set_eq:nn { #1 } { ltxtalk-#1 }
  }

\NewDocumentCommand \setltxtalkcolors { m }
  {
    \keys_set:nn { ltxtalk / colors } { #1 }
  }

% Standardfarben
\color_set:nnn { ltxtalk-primary } { HTML } { 2B3A67 }
\color_set:nnn { ltxtalk-secondary } { HTML } { 496A81 }
\color_set:nnn { ltxtalk-accent } { HTML } { 66999B }
\color_set:nnn { ltxtalk-text } { HTML } { 333333 }
\color_set:nnn { ltxtalk-bg } { HTML } { FFFFFF }
\__ltxtalk_set_semantic_color:nn { structure } { {HTML}{2B3A67} }
\__ltxtalk_set_semantic_color:nn { alert }     { {HTML}{B00020} }
\__ltxtalk_set_semantic_color:nn { example }   { {HTML}{166534} }

% ==========================================
% INNER THEME: BLOCKS AND LIST MARKERS
% ==========================================

\tl_new:N \l__ltxtalk_block_background_tl
\tl_new:N \l__ltxtalk_block_rule_color_tl
\tl_new:N \l__ltxtalk_block_before_tl
\tl_new:N \l__ltxtalk_block_after_tl
\dim_new:N \l__ltxtalk_block_rule_dim
\dim_new:N \l__ltxtalk_block_padding_dim
\dim_new:N \l__ltxtalk_block_corner_radius_dim
\dim_new:N \l__ltxtalk_block_shadow_offset_dim
\bool_new:N \l__ltxtalk_block_shadow_bool

\NewTemplateType { lttheme-block } { 1 }
\DeclareTemplateInterface { lttheme-block } { standard } { 1 }
  {
    background    : tokenlist = none,
    rule-color    : tokenlist = structure,
    rule-thickness: length = 0pt,
    padding       : length = 0pt,
    corner-radius : length = 0pt,
    shadow        : boolean = false,
    shadow-offset : length = 2pt,
    before        : tokenlist =,
    after         : tokenlist =
  }
\DeclareTemplateCode { lttheme-block } { standard } { 1 }
  {
    background = \l__ltxtalk_block_background_tl,
    rule-color = \l__ltxtalk_block_rule_color_tl,
    rule-thickness = \l__ltxtalk_block_rule_dim,
    padding = \l__ltxtalk_block_padding_dim,
    corner-radius = \l__ltxtalk_block_corner_radius_dim,
    shadow = \l__ltxtalk_block_shadow_bool,
    shadow-offset = \l__ltxtalk_block_shadow_offset_dim,
    before = \l__ltxtalk_block_before_tl,
    after = \l__ltxtalk_block_after_tl
  }
  { \__ltxtalk_typeset_block:n {#1} }

\cs_new_protected:Npn \__ltxtalk_typeset_block:n #1
  {
    \group_begin:
      \l__ltxtalk_block_before_tl
      \dim_compare:nNnTF { \l__ltxtalk_block_rule_dim } > { 0pt }
        {
          \setlength\fboxsep{\l__ltxtalk_block_padding_dim}
          \setlength\fboxrule{\l__ltxtalk_block_rule_dim}
          \noindent\fcolorbox{\l__ltxtalk_block_rule_color_tl}
            {\str_if_eq:VnTF \l__ltxtalk_block_background_tl {none}
               {ltxtalk-bg}{\l__ltxtalk_block_background_tl}}
            {\begin{minipage}{\dim_eval:n
                {\linewidth-2\fboxsep-2\fboxrule}}#1\end{minipage}}
        }
        { #1 }
      \l__ltxtalk_block_after_tl
    \group_end:
  }

\clist_map_inline:nn { block, alertblock, exampleblock }
  { \DeclareInstance { lttheme-block } {#1} { standard } { } }

\NewDocumentCommand \styleltxtalkblock { m m }
  { \EditInstance { lttheme-block } {#1} {#2} }

\cs_new_protected:Npn \__ltxtalk_use_original_block:nnn #1#2#3
  {
    \IfNoValueTF {#1}
      { \begin{ltxtalk-original-block}{#2}#3\end{ltxtalk-original-block} }
      { \begin{ltxtalk-original-block}<#1>{#2}#3\end{ltxtalk-original-block} }
  }

\NewEnvironmentCopy { ltxtalk-original-block } { block }
\RenewDocumentEnvironment { block } { d <> m +b }
  { \UseInstance { lttheme-block } { block }
      {\__ltxtalk_use_original_block:nnn {#1}{#2}{#3}} }
  { }

% Only provide the two variants while the class does not provide its own.
\cs_if_exist:cF { alertblock }
  {
    \NewDocumentEnvironment { alertblock } { d <> m +b }
      { \UseInstance { lttheme-block } { alertblock }
          {\group_begin:\color_set_eq:nn {structure}{alert}
           \__ltxtalk_use_original_block:nnn {#1}{#2}{#3}\group_end:} }
      { }
  }
\cs_if_exist:cF { exampleblock }
  {
    \NewDocumentEnvironment { exampleblock } { d <> m +b }
      { \UseInstance { lttheme-block } { exampleblock }
          {\group_begin:\color_set_eq:nn {structure}{example}
           \__ltxtalk_use_original_block:nnn {#1}{#2}{#3}\group_end:} }
      { }
  }

\keys_define:nn { ltxtalk / itemize }
  {
    level-1 .code:n = \cs_set:Npn \labelitemi {#1},
    level-2 .code:n = \cs_set:Npn \labelitemii {#1},
    level-3 .code:n = \cs_set:Npn \labelitemiii {#1},
    level-4 .code:n = \cs_set:Npn \labelitemiv {#1}
  }
\NewDocumentCommand \setltxtalkitemize { m }
  { \keys_set:nn { ltxtalk / itemize } {#1} }

% ==========================================
% NAMED-SLOT-SYSTEM
% ==========================================

\prop_new:N \g__ltxtalk_slot_types_prop
\prop_new:N \g__ltxtalk_slot_positions_prop
\prop_new:N \g__ltxtalk_slot_styles_prop
\prop_new:N \g__ltxtalk_slot_areas_prop
\prop_new:N \g__ltxtalk_slot_accessibility_prop
\prop_new:N \g__ltxtalk_area_grid_columns_prop
\prop_new:N \g__ltxtalk_area_grid_rows_prop
\prop_new:N \g__ltxtalk_slot_grid_positions_prop

\tl_new:N \l__ltxtalk_current_slot_name_tl
\tl_new:N \l__ltxtalk_slot_type_tl
\tl_new:N \l__ltxtalk_slot_content_tl
\tl_new:N \l__ltxtalk_slot_style_tl
\tl_new:N \l__ltxtalk_slot_area_tl
\tl_new:N \l__ltxtalk_slot_accessibility_tl

\NewDocumentCommand \defineltxtalkslot { m m m m }
  {
    % \begin{itemize}
    % \item \#1: Slot-Name (eindeutig)
    % \item \#2: Bereich (header/left/content/right/footer)
    % \item \#3: Slot-Typ (frametitle, section, logo, etc.)
    % \item \#4: Accessibility-Optionen (key-value)
    % \end{itemize}
    \prop_gput:Nnn \g__ltxtalk_slot_areas_prop { #1 } { #2 }
    \prop_gput:Nnn \g__ltxtalk_slot_types_prop { #1 } { #3 }
    \prop_gput:Nnn \g__ltxtalk_slot_accessibility_prop { #1 } { #4 }
    
    % Automatische Rollenerkennung basierend auf Slot-Typ
    \str_case:nnF { #3 }
      {
        { frametitle }    { \keys_set:nn { ltxtalk / a11y } { role = heading, heading-level = 2 } }
        { framesubtitle } { \keys_set:nn { ltxtalk / a11y } { role = heading, heading-level = 3 } }
        { title }         { \keys_set:nn { ltxtalk / a11y } { role = heading, heading-level = 1 } }
        { section }       { \keys_set:nn { ltxtalk / a11y } { role = heading, heading-level = 2 } }
        { subsection }    { \keys_set:nn { ltxtalk / a11y } { role = heading, heading-level = 3 } }
        { pagenumber }    { \keys_set:nn { ltxtalk / a11y } { role = navigation } }
        { logo }          { \keys_set:nn { ltxtalk / a11y } { role = artifact, alt-text = Logo } }
      }
      {
        \keys_set:nn { ltxtalk / a11y } { role = artifact }
      }
  }

\NewDocumentCommand \positionltxtalkslot { m m m m m }
  {
    % 
    % \begin{itemize}
    %  \item \#1: Slot-Name
    %  \item \#2: x-Position (Dimension oder 'left'/'center'/'right')
    %  \item \#3: y-Position (Dimension oder 'top'/'center'/'bottom')
    %  \item \#4: Breite (Dimension)
    %  \item \#5: Höhe (Dimension)
    %\end{itemize}
    \prop_gput:Nnn \g__ltxtalk_slot_positions_prop { #1 } 
      { { #2 } { #3 } { #4 } { #5 } }
  }

\NewDocumentCommand \DeclareLTXTalkAreaGrid { m m m }
  {
    \prop_gput:Nnn \g__ltxtalk_area_grid_columns_prop {#1} {#2}
    \prop_gput:Nnn \g__ltxtalk_area_grid_rows_prop {#1} {#3}
  }

\NewDocumentCommand \PositionLTXTalkSlotInGrid { m m m m m }
  {\prop_gput:Nnn \g__ltxtalk_slot_grid_positions_prop {#1}
    {{#2}{#3}{#4}{#5}}}

\NewDocumentCommand \SetLTXTalkAreaGridRows { m m }
  {\prop_gput:Nnn \g__ltxtalk_area_grid_rows_prop {#1}{#2}}

% ==========================================
% SLOT-FORMATIERUNG
% ==========================================

\tl_new:N \l__ltxtalk_slot_font_tl
\tl_new:N \l__ltxtalk_slot_color_tl
\tl_new:N \l__ltxtalk_slot_bg_tl
\tl_new:N \l__ltxtalk_slot_align_tl
\tl_new:N \l__ltxtalk_slot_frame_tl
\tl_new:N \l__ltxtalk_slot_rule_tl
\tl_new:N \l__ltxtalk_slot_rotate_tl
\tl_new:N \l__ltxtalk_slot_opacity_tl

\dim_new:N \l__ltxtalk_slot_padding_dim
\dim_new:N \l__ltxtalk_slot_rule_thick_dim

\bool_new:N \l__ltxtalk_slot_shadow_bool
\bool_new:N \l__ltxtalk_slot_uppercase_bool
\bool_new:N \l__ltxtalk_slot_lowercase_bool

\keys_define:nn { ltxtalk / slot-style }
  {
    font        .tl_set:N = \l__ltxtalk_slot_font_tl,
    font        .initial:n = \normalfont,
    color       .tl_set:N = \l__ltxtalk_slot_color_tl,
    color       .initial:n = ltxtalk-text,
    background  .tl_set:N = \l__ltxtalk_slot_bg_tl,
    background  .initial:n = none,
    align       .tl_set:N = \l__ltxtalk_slot_align_tl,
    align       .initial:n = center,
    frame       .tl_set:N = \l__ltxtalk_slot_frame_tl,
    frame       .initial:n = none,
    padding     .dim_set:N = \l__ltxtalk_slot_padding_dim,
    padding     .initial:n = 0pt,
    rule-color  .tl_set:N = \l__ltxtalk_slot_rule_tl,
    rule-color  .initial:n = none,
    rule-thick  .dim_set:N = \l__ltxtalk_slot_rule_thick_dim,
    rule-thick  .initial:n = 0.4pt,
    
    rotate      .tl_set:N = \l__ltxtalk_slot_rotate_tl,
    rotate      .initial:n = 0,
    opacity     .tl_set:N = \l__ltxtalk_slot_opacity_tl,
    opacity     .initial:n = 1,
    shadow      .bool_set:N = \l__ltxtalk_slot_shadow_bool,
    shadow      .initial:n = false,
    uppercase   .bool_set:N = \l__ltxtalk_slot_uppercase_bool,
    uppercase   .initial:n = false,
    lowercase   .bool_set:N = \l__ltxtalk_slot_lowercase_bool,
    lowercase   .initial:n = false,
  }

\NewDocumentCommand \styleltxtalkslot { m m }
  {
    % \begin{itemize}
    %  \item \#1: Slot-Name
    %  \item \#2: Formatierungsoptionen (key-value)
    % \end{itemize} 
    \prop_gput:Nnn \g__ltxtalk_slot_styles_prop { #1 } { #2 }
  }

% ==========================================
% BEREICHSDEKORATIONEN
% ==========================================

\prop_new:N \g__ltxtalk_area_decorations_prop

\tl_new:N \l__ltxtalk_area_deco_tl
\tl_new:N \l__ltxtalk_area_frame_tl
\tl_new:N \l__ltxtalk_area_rule_color_tl
\tl_new:N \l__ltxtalk_area_rule_style_tl
\tl_new:N \l__ltxtalk_area_background_tl
\tl_new:N \l__ltxtalk_area_bg_opacity_tl
\tl_new:N \l__ltxtalk_area_shadow_color_tl
\tl_new:N \l__ltxtalk_artifact_type_tl

\dim_new:N \l__ltxtalk_area_rule_thickness_dim
\dim_new:N \l__ltxtalk_area_shadow_offset_dim
\dim_new:N \l__ltxtalk_area_padding_dim
\dim_new:N \l__ltxtalk_area_margin_top_dim
\dim_new:N \l__ltxtalk_area_margin_bottom_dim

\bool_new:N \l__ltxtalk_area_shadow_bool

\NewDocumentCommand \decorateltxtalkarea { m m }
  {
    % 
    % \begin{itemize}
    %  \item \#1: Bereich (header/left/right/footer/content)
    %  \item \#2: Dekorationsoptionen (key-value)
    % \end{itemize}
    \prop_gput:Nnn \g__ltxtalk_area_decorations_prop { #1 } { #2 }
  }

\keys_define:nn { ltxtalk / area-decoration }
  {
    frame           .choice:,
    frame / none    .code:n = \tl_set:Nn \l__ltxtalk_area_frame_tl { none },
    frame / box     .code:n = \tl_set:Nn \l__ltxtalk_area_frame_tl { box },
    frame / top     .code:n = \tl_set:Nn \l__ltxtalk_area_frame_tl { top },
    frame / bottom  .code:n = \tl_set:Nn \l__ltxtalk_area_frame_tl { bottom },
    frame / left    .code:n = \tl_set:Nn \l__ltxtalk_area_frame_tl { left },
    frame / right   .code:n = \tl_set:Nn \l__ltxtalk_area_frame_tl { right },
    frame / topbottom .code:n = \tl_set:Nn \l__ltxtalk_area_frame_tl { topbottom },
    frame .initial:n = none,
    
    rule-color      .tl_set:N = \l__ltxtalk_area_rule_color_tl,
    rule-color      .initial:n = ltxtalk-primary,
    rule-thickness  .dim_set:N = \l__ltxtalk_area_rule_thickness_dim,
    rule-thickness  .initial:n = 1pt,
    rule-style      .tl_set:N = \l__ltxtalk_area_rule_style_tl,
    rule-style      .initial:n = solid,
    
    background      .tl_set:N = \l__ltxtalk_area_background_tl,
    background      .initial:n = none,
    background-opacity .tl_set:N = \l__ltxtalk_area_bg_opacity_tl,
    background-opacity .initial:n = 1,
    
    shadow          .bool_set:N = \l__ltxtalk_area_shadow_bool,
    shadow          .initial:n = false,
    shadow-color    .tl_set:N = \l__ltxtalk_area_shadow_color_tl,
    shadow-color    .initial:n = black!30,
    shadow-offset   .dim_set:N = \l__ltxtalk_area_shadow_offset_dim,
    shadow-offset   .initial:n = 3pt,
    
    padding         .dim_set:N = \l__ltxtalk_area_padding_dim,
    padding         .initial:n = 0pt,
    margin-top      .dim_set:N = \l__ltxtalk_area_margin_top_dim,
    margin-top      .initial:n = 0pt,
    margin-bottom   .dim_set:N = \l__ltxtalk_area_margin_bottom_dim,
    margin-bottom   .initial:n = 0pt,
    
    artifact-type   .choice:,
    artifact-type / decorative  .code:n = \tl_set:Nn \l__ltxtalk_artifact_type_tl { decorative },
    artifact-type / layout      .code:n = \tl_set:Nn \l__ltxtalk_artifact_type_tl { layout },
    artifact-type / background  .code:n = \tl_set:Nn \l__ltxtalk_artifact_type_tl { background },
    artifact-type .initial:n = decorative,
  }

% ==========================================
% TITEL-PLATZIERUNG
% ==========================================

\tl_new:N \g__ltxtalk_frametitle_target_tl
\tl_gset:Nn \g__ltxtalk_frametitle_target_tl { content }
\tl_new:N \g__ltxtalk_section_target_tl
\tl_gset:Nn \g__ltxtalk_section_target_tl { none }

\NewDocumentCommand \setltxtalktitleplacement { m }
  {
    \keys_set:nn { ltxtalk / title-placement } { #1 }
  }

\keys_define:nn { ltxtalk / title-placement }
  {
    frametitle .choice:,
    frametitle / content .code:n = 
      { \tl_gset:Nn \g__ltxtalk_frametitle_target_tl { content } },
    frametitle / header .code:n = 
      { \tl_gset:Nn \g__ltxtalk_frametitle_target_tl { header } },
    frametitle / both .code:n = 
      { \tl_gset:Nn \g__ltxtalk_frametitle_target_tl { both } },
    frametitle .initial:n = content,
    
    section .choice:,
    section / header .code:n = 
      { \tl_gset:Nn \g__ltxtalk_section_target_tl { header } },
    section / none .code:n = 
      { \tl_gset:Nn \g__ltxtalk_section_target_tl { none } },
    section .initial:n = none,
  }

% ==========================================
% ACCESSIBILITY
% ==========================================

\tl_new:N \l__ltxtalk_slot_role_tl
\tl_new:N \l__ltxtalk_alt_text_tl
\tl_new:N \l__ltxtalk_continuation_text_tl
\tl_new:N \l__ltxtalk_heading_hash_tl
\tl_new:N \l__ltxtalk_heading_tag_tl
\tl_new:N \l__ltxtalk_clean_title_tl
\tl_new:N \l__ltxtalk_previous_context_tl

\int_new:N \l__ltxtalk_heading_level_int

\keys_define:nn { ltxtalk / a11y }
  {
    role .choice:,
    role / heading       .code:n = \tl_set:Nn \l__ltxtalk_slot_role_tl { heading },
    role / heading-1     .code:n = \tl_set:Nn \l__ltxtalk_slot_role_tl { H1 },
    role / heading-2     .code:n = \tl_set:Nn \l__ltxtalk_slot_role_tl { H2 },
    role / heading-3     .code:n = \tl_set:Nn \l__ltxtalk_slot_role_tl { H3 },
    role / artifact      .code:n = \tl_set:Nn \l__ltxtalk_slot_role_tl { artifact },
    role / content       .code:n = \tl_set:Nn \l__ltxtalk_slot_role_tl { content },
    role / navigation    .code:n = \tl_set:Nn \l__ltxtalk_slot_role_tl { navigation },
    role / note          .code:n = \tl_set:Nn \l__ltxtalk_slot_role_tl { note },
    role .initial:n = artifact,
    
    heading-level .int_set:N = \l__ltxtalk_heading_level_int,
    heading-level .initial:n = 1,
    
    alt-text .tl_set:N = \l__ltxtalk_alt_text_tl,
    alt-text .initial:n = {},
    
    continuation-text .tl_set:N = \l__ltxtalk_continuation_text_tl,
    continuation-text .initial:n = { (Fortsetzung) },
    continuation-mode .choice:,
    continuation-mode / explicit .code:n = 
      { \bool_set_true:N \l__ltxtalk_explicit_cont_bool },
    continuation-mode / implicit .code:n = 
      { \bool_set_false:N \l__ltxtalk_explicit_cont_bool },
    continuation-mode .initial:n = implicit,
  }

\NewDocumentCommand \setltxtalkheadingbehavior { m }
  {
    \keys_set:nn { ltxtalk / heading } { #1 }
  }

\keys_define:nn { ltxtalk / heading }
  {
    auto-section-title .bool_set:N = \l__ltxtalk_auto_section_title_bool,
    auto-section-title .initial:n = true,
    detect-continuation .bool_set:N = \l__ltxtalk_detect_cont_bool,
    detect-continuation .initial:n = true,
    header-as-navigation .bool_set:N = \l__ltxtalk_header_nav_bool,
    header-as-navigation .initial:n = true,
  }

% Heading-Tracking
\prop_new:N \g__ltxtalk_heading_context_prop
\seq_new:N \g__ltxtalk_heading_history_seq

% ==========================================
% FRAME-INFORMATIONEN SAMMELN
% ==========================================

\tl_new:N \g__ltxtalk_current_frametitle_tl
\tl_new:N \g__ltxtalk_current_framesubtitle_tl
\tl_new:N \g__ltxtalk_current_section_tl
\tl_new:N \g__ltxtalk_current_subsection_tl

\cs_new_protected:Npn \__ltxtalk_capture_frame_info:
  {
    \cs_if_exist:NT \insertframetitle
      {
        \tl_if_empty:NF \insertframetitle
          {
            \tl_gset:Nx \g__ltxtalk_current_frametitle_tl { \insertframetitle }
          }
      }
    
    \cs_if_exist:NT \insertframesubtitle
      {
        \tl_if_empty:NF \insertframesubtitle
          {
            \tl_gset:Nx \g__ltxtalk_current_framesubtitle_tl { \insertframesubtitle }
          }
      }
    
    \cs_if_exist:NT \insertsection
      {
        \tl_gset:Nx \g__ltxtalk_current_section_tl { \insertsection }
      }
    
    \cs_if_exist:NT \insertsubsection
      {
        \tl_gset:Nx \g__ltxtalk_current_subsection_tl { \insertsubsection }
      }
  }

% ==========================================
% HEADING-QUELLENERKENNUNG
% ==========================================

\cs_new_protected:Npn \__ltxtalk_detect_heading_source:
  {
    \bool_if:NT \l__ltxtalk_auto_section_title_bool
      {
        \tl_if_empty:NF \g__ltxtalk_current_frametitle_tl
          {
            \str_if_eq:eeTF 
              { \g__ltxtalk_current_frametitle_tl } 
              { \g__ltxtalk_current_section_tl }
              {
                \keys_set:nn { ltxtalk / a11y } {
                  role = heading,
                  heading-level = 2,
                }
              }
              {
                \str_if_eq:eeTF 
                  { \g__ltxtalk_current_frametitle_tl } 
                  { \g__ltxtalk_current_subsection_tl }
                  {
                    \keys_set:nn { ltxtalk / a11y } {
                      role = heading,
                      heading-level = 3,
                    }
                  }
                  {}
              }
          }
      }
  }

% ==========================================
% FORTSETZUNGSFRAME-ERKENNUNG
% ==========================================

\cs_new_protected:Npn \__ltxtalk_handle_continuation_frame:n #1
  {
    \bool_if:NT \l__ltxtalk_detect_cont_bool
      {
        \tl_if_in:nnTF { #1 } { \l__ltxtalk_continuation_text_tl }
          {
            \__ltxtalk_mark_as_continuation:n { #1 }
          }
          {
            \__ltxtalk_check_implicit_continuation:n { #1 }
          }
      }
  }

\cs_new_protected:Npn \__ltxtalk_mark_as_continuation:n #1
  {
    \tl_set:Nx \l__ltxtalk_clean_title_tl {
      \tl_trim_spaces:n {
        \str_replace_all:nnn { #1 } { \l__ltxtalk_continuation_text_tl } {}
      }
    }
    
    \prop_gput:Nnn \g__ltxtalk_heading_context_prop 
      { heading- \l__ltxtalk_clean_title_tl } 
      { continued }
  }

\cs_new_protected:Npn \__ltxtalk_check_implicit_continuation:n #1
  {
    \bool_if:NF \l__ltxtalk_explicit_cont_bool
      {
        \prop_get:NnNTF \g__ltxtalk_heading_context_prop 
          { heading-#1 } 
          \l__ltxtalk_previous_context_tl
          {
            \str_if_eq:VnT \l__ltxtalk_previous_context_tl { frame }
              {
                \__ltxtalk_mark_as_continuation:n { #1~(cont.) }
              }
          }
          {}
      }
  }

% ==========================================
% DOPPELTE TITEL VERMEIDEN
% ==========================================

\cs_new_protected:Npn \__ltxtalk_manage_frametitle_placement:
  {
    \bool_set_false:N \l__ltxtalk_frametitle_in_header_bool
    
    \prop_map_inline:Nn \g__ltxtalk_slot_types_prop
      {
        \str_if_eq:nnTF { ##2 } { frametitle }
          {
            \prop_get:NnN \g__ltxtalk_slot_areas_prop { ##1 } \l_tmpa_tl
            \str_if_eq:VnT \l_tmpa_tl { header }
              {
                \bool_set_true:N \l__ltxtalk_frametitle_in_header_bool
                \prop_map_break:
              }
          }
          {}
      }
  }

% ==========================================
% SLOT-CONTENT-RESOLVER
% ==========================================

\tl_new:N \l__ltxtalk_metadata_lines_tl
\tl_new:N \l__ltxtalk_metadata_linebreak_tl
\keys_define:nn {ltxtalk/metadata-use}
  {
    lines .choice:,
    lines / all .code:n = \tl_set:Nn \l__ltxtalk_metadata_lines_tl {all},
    lines / 1 .code:n = \tl_set:Nn \l__ltxtalk_metadata_lines_tl {1},
    lines .initial:n = all,
    linebreak .tl_set:N = \l__ltxtalk_metadata_linebreak_tl,
    linebreak .initial:n = {~}
  }

\cs_new:Npn \__ltxtalk_metadata:n #1
  {
    \str_case:nnF {#1}
      {
        {title}{\@title}
        {short-title}{\@shorttitle}
        {subtitle}{\@subtitle}
        {short-subtitle}{\@shortsubtitle}
        {author}{\@author}
        {short-author}{\@shortauthor}
        {institute}{\@institute}
        {short-institute}{\@shortinstitute}
        {date}{\@date}
        {short-date}{\@shortdate}
        {section}{\insertsection}
        {subsection}{\insertsubsection}
        {frametitle}{\tl_use:N \g__talk_frame_title_tl}
        {framesubtitle}{\tl_use:N \g__talk_frame_subtitle_tl}
        {framenumber}{\@framenumber}
        {totalframes}{\@totalframes}
      }
      { }
  }

\NewExpandableDocumentCommand \LTXTalkMetadata {m}
  {\__ltxtalk_metadata:n {#1}}

\NewDocumentCommand \UseLTXTalkMetadata {O{} m}
  {
    \group_begin:
      \keys_set:nn {ltxtalk/metadata-use}{#1}
      \tl_set:Ne \l_tmpa_tl {\__ltxtalk_metadata:n {#2}}
      \str_if_eq:VnT \l__ltxtalk_metadata_lines_tl {1}
        {\tl_replace_all:NnV \l_tmpa_tl {\\} \l__ltxtalk_metadata_linebreak_tl}
      \tl_use:N \l_tmpa_tl
    \group_end:
  }

\cs_new_protected:Npn \__ltxtalk_resolve_slot_content:Nn #1 #2
  {
    % 
    % \begin{itemize}
    %  \item \#1: Slot-Typ
    %  \item \#2: Ziel-Tokenlist
    % \end{itemize}
    
    \str_case:nn { #1 }
      {
        { title }      { \tl_set:Nn #2 { \LTXTalkMetadata{short-title} } }
        { subtitle }   { \tl_set:Nn #2 { \LTXTalkMetadata{short-subtitle} } }
        { author }     { \tl_set:Nn #2 { \LTXTalkMetadata{short-author} } }
        { date }       { \tl_set:Nn #2 { \LTXTalkMetadata{date} } }
        { institute }  { \tl_set:Nn #2 { \LTXTalkMetadata{short-institute} } }
        { section }    { \tl_set:Nn #2 { \g__ltxtalk_current_section_tl } }
        { subsection } { \tl_set:Nn #2 { \g__ltxtalk_current_subsection_tl } }
        { frametitle } { \tl_set:Nn #2 { \g__ltxtalk_current_frametitle_tl } }
        { framesubtitle } { \tl_set:Nn #2 { \g__ltxtalk_current_framesubtitle_tl } }
        { pagenumber } { \tl_set:Nn #2 { \insertframenumber } }
        { logo }       { \tl_set:Nn #2 { \includegraphics[height=1cm]{\insertlogo} } }
      }
  }

\cs_generate_variant:Nn \__ltxtalk_resolve_slot_content:Nn { NV }

% ==========================================
% SLOT-RENDERING
% ==========================================

\cs_new_protected:Npn \__ltxtalk_render_named_slot:n #1
  {
    \prop_get:NnNTF \g__ltxtalk_slot_types_prop { #1 } \l__ltxtalk_slot_type_tl
      {
        \__ltxtalk_resolve_slot_content:NV \l__ltxtalk_slot_type_tl \l__ltxtalk_slot_content_tl
        
        % Formatierung anwenden
        \prop_get:NnNTF \g__ltxtalk_slot_styles_prop { #1 } \l__ltxtalk_slot_style_tl
          {
            \keys_set:nV { ltxtalk / slot-style } \l__ltxtalk_slot_style_tl
          }
          {
            \__ltxtalk_default_slot_style:V \l__ltxtalk_slot_type_tl
          }
        
        % Accessibility-Tagging
        \__ltxtalk_tag_slot_content:n { \l__ltxtalk_slot_content_tl }
      }
      {}
  }

% ==========================================
% DEFAULT-FORMATIERUNG
% ==========================================

\cs_new_protected:Npn \__ltxtalk_default_slot_style:n #1
  {
    \str_case:nn { #1 }
      {
        { title }      { 
          \keys_set:nn { ltxtalk / slot-style } {
            font = \Large\bfseries\sffamily,
            color = ltxtalk-primary,
          }
        }
        { frametitle } { 
          \keys_set:nn { ltxtalk / slot-style } {
            font = \Large\bfseries,
            color = ltxtalk-primary,
            frame = bottom,
            rule-color = ltxtalk-primary,
            rule-thick = 1pt,
          }
        }
        { author }     { 
          \keys_set:nn { ltxtalk / slot-style } {
            font = \small\itshape,
            color = ltxtalk-secondary,
          }
        }
        { date }       { 
          \keys_set:nn { ltxtalk / slot-style } {
            font = \small,
            color = ltxtalk-secondary!70,
          }
        }
        { section }    { 
          \keys_set:nn { ltxtalk / slot-style } {
            font = \normalsize\bfseries,
            color = ltxtalk-accent,
          }
        }
        { subsection } { 
          \keys_set:nn { ltxtalk / slot-style } {
            font = \small\bfseries,
            color = ltxtalk-accent!80,
          }
        }
        { pagenumber } { 
          \keys_set:nn { ltxtalk / slot-style } {
            font = \tiny\sffamily,
            color = ltxtalk-secondary!50,
          }
        }
        { framesubtitle } { 
          \keys_set:nn { ltxtalk / slot-style } {
            font = \large\itshape,
            color = ltxtalk-secondary,
          }
        }
      }
  }

\cs_generate_variant:Nn \__ltxtalk_default_slot_style:n { V }

% ==========================================
% ACCESSIBILITY-TAGGING
% ==========================================

\cs_new_protected:Npn \__ltxtalk_tag_slot_content:n #1
  {
    \prop_get:NnNTF \g__ltxtalk_slot_accessibility_prop
      { \l__ltxtalk_current_slot_name_tl } 
      \l__ltxtalk_slot_accessibility_tl
      {
        \keys_set:nV { ltxtalk / a11y } \l__ltxtalk_slot_accessibility_tl
      }
      {
        \keys_set:nn { ltxtalk / a11y } { role = artifact }
      }
    
    \str_case:Vn \l__ltxtalk_slot_role_tl
      {
        { heading } {
          \__ltxtalk_track_heading_appearance:nnn 
            { \l__ltxtalk_heading_level_int } 
            { #1 }
            { frame }
        }
        { H1 } {
          \__ltxtalk_track_heading_appearance:nnn { 1 } { #1 } { frame }
        }
        { H2 } {
          \__ltxtalk_track_heading_appearance:nnn { 2 } { #1 } { frame }
        }
        { H3 } {
          \__ltxtalk_track_heading_appearance:nnn { 3 } { #1 } { frame }
        }
        { artifact } {
          \__ltxtalk_render_artifact:n { #1 }
        }
        { navigation } {
          \__ltxtalk_render_navigation:n { #1 }
        }
        { content } {
          \__ltxtalk_render_content:n { #1 }
        }
        { note } {
          \__ltxtalk_render_note:n { #1 }
        }
      }
  }

\cs_new_protected:Npn \__ltxtalk_track_heading_appearance:nnn #1 #2 #3
  {
    \int_gincr:N \g__ltxtalk_heading_id_int
    
    \tl_set:Nx \l__ltxtalk_heading_hash_tl {
      H#1- \tl_to_str:n { #2 } - \tl_to_str:n { #3 }
    }
    
    \prop_get:NnNTF \g__ltxtalk_heading_context_prop 
      { \l__ltxtalk_heading_hash_tl } 
      \l__ltxtalk_previous_context_tl
      {
        % Wiederholtes Heading -> Artifact
        \__ltxtalk_render_artifact:n { #2 }
      }
      {
        % Erstmaliges Heading -> Semantisch
        \prop_gput:Nnn \g__ltxtalk_heading_context_prop 
          { \l__ltxtalk_heading_hash_tl } 
          { #3- \int_use:N \g__ltxtalk_heading_id_int }
        
        \__ltxtalk_render_heading:nn { #1 } { #2 }
      }
  }

\cs_new_protected:Npn \__ltxtalk_render_heading:nn #1 #2
  {
    \tl_set:Nx \l__ltxtalk_heading_tag_tl { H#1 }
    % Hier würde echtes PDF-Tagging stattfinden
    #2
  }

\cs_new_protected:Npn \__ltxtalk_render_artifact:n #1
  {
    % Hier als Artifact taggen
    #1
  }

\cs_new_protected:Npn \__ltxtalk_render_navigation:n #1
  {
    % Hier als Navigation taggen
    #1
  }

\cs_new_protected:Npn \__ltxtalk_render_content:n #1
  {
    % Hier als Content taggen
    #1
  }

\cs_new_protected:Npn \__ltxtalk_render_note:n #1
  {
    % Hier als Note taggen
    #1
  }

% ==========================================
% BEREICHS-RENDERING
% ==========================================

\cs_new_protected:Npn \__ltxtalk_render_area_with_decorations:nn #1 #2
  {
    % \begin{itemize}
    %  \item \#1: Bereichsname
    %  \item \#2: Inhalt (Slots)
    % \end{itemize} 
    
    \prop_get:NnNTF \g__ltxtalk_area_decorations_prop { #1 } \l__ltxtalk_area_deco_tl
      {
        \keys_set:nV { ltxtalk / area-decoration } \l__ltxtalk_area_deco_tl
      }
      {
        \tl_set:Nn \l__ltxtalk_area_frame_tl { none }
        \tl_set:Nn \l__ltxtalk_area_background_tl { none }
      }
    
    % Hintergrund (als Artifact)
    \str_if_eq:VnF \l__ltxtalk_area_background_tl { none }
      {
        \__ltxtalk_render_artifact:n {
          \color{ \l__ltxtalk_area_background_tl }
          \rule{ \l__ltxtalk_area_width_dim }{ \l__ltxtalk_area_height_dim }
        }
      }
    
    % Inhalt rendern
    #2
    
    % Rahmen (als Artifact)
    \str_if_eq:VnF \l__ltxtalk_area_frame_tl { none }
      {
        \__ltxtalk_render_artifact:n {
          \__ltxtalk_draw_area_frame:
        }
      }
  }

\cs_new_protected:Npn \__ltxtalk_draw_area_frame:
  {
    \color{ \l__ltxtalk_area_rule_color_tl }
    \dim_set:Nn \l_tmpa_dim { \l__ltxtalk_area_rule_thickness_dim }
    
    \str_case:Vn \l__ltxtalk_area_frame_tl
      {
        { bottom } {
          \rule{ \l__ltxtalk_area_width_dim }{ \l_tmpa_dim }
        }
        { top } {
          \rule{ \l__ltxtalk_area_width_dim }{ \l_tmpa_dim }
        }
        { topbottom } {
          \rule{ \l__ltxtalk_area_width_dim }{ \l_tmpa_dim }
          \rule{ \l__ltxtalk_area_width_dim }{ \l_tmpa_dim }
        }
      }
  }

% ==========================================
% SEITENAUFBAU
% ==========================================

\cs_new_protected:Npn \__ltxtalk_build_themed_page:
  {
    \str_if_eq:VnF \g__ltxtalk_active_theme_tl { default }
      {
        \vcoffin_set:Nnn \g__ltxtalk_page_coffin 
          { \paperwidth } { \paperheight }
          {
            \color{ltxtalk-bg}
            \rule{\paperwidth}{\paperheight}
          }
        
        % Header rendern
        \bool_if:NT \l__ltxtalk_show_header_bool
          {
            \__ltxtalk_render_area_with_decorations:nn { header }
              {
                % Alle Header-Slots rendern
                \prop_map_inline:Nn \g__ltxtalk_slot_areas_prop
                  {
                    \str_if_eq:nnT { ##2 } { header }
                      {
                        \tl_set:Nn \l__ltxtalk_current_slot_name_tl { ##1 }
                        \__ltxtalk_render_named_slot:n { ##1 }
                      }
                  }
              }
          }
        
        % Linke Margin
        \bool_if:NT \l__ltxtalk_show_left_margin_bool
          {
            \__ltxtalk_render_area_with_decorations:nn { left }
              {
                \prop_map_inline:Nn \g__ltxtalk_slot_areas_prop
                  {
                    \str_if_eq:nnT { ##2 } { left }
                      {
                        \tl_set:Nn \l__ltxtalk_current_slot_name_tl { ##1 }
                        \__ltxtalk_render_named_slot:n { ##1 }
                      }
                  }
              }
          }
        
        % Content
        \prop_map_inline:Nn \g__ltxtalk_slot_areas_prop
          {
            \str_if_eq:nnT { ##2 } { content }
              {
                \tl_set:Nn \l__ltxtalk_current_slot_name_tl { ##1 }
                \__ltxtalk_render_named_slot:n { ##1 }
              }
          }
        
        % Rechte Margin
        \bool_if:NT \l__ltxtalk_show_right_margin_bool
          {
            \__ltxtalk_render_area_with_decorations:nn { right }
              {
                \prop_map_inline:Nn \g__ltxtalk_slot_areas_prop
                  {
                    \str_if_eq:nnT { ##2 } { right }
                      {
                        \tl_set:Nn \l__ltxtalk_current_slot_name_tl { ##1 }
                        \__ltxtalk_render_named_slot:n { ##1 }
                      }
                  }
              }
          }
        
        % Footer
        \bool_if:NT \l__ltxtalk_show_footer_bool
          {
            \__ltxtalk_render_area_with_decorations:nn { footer }
              {
                \prop_map_inline:Nn \g__ltxtalk_slot_areas_prop
                  {
                    \str_if_eq:nnT { ##2 } { footer }
                      {
                        \tl_set:Nn \l__ltxtalk_current_slot_name_tl { ##1 }
                        \__ltxtalk_render_named_slot:n { ##1 }
                      }
                  }
              }
          }
        
        \coffin_typeset:Nnnnn \g__ltxtalk_page_coffin 
          { l } { t } { 0pt } { 0pt }
      }
  }

% ==========================================
% INTEGRATION IN DIE LTX-TALK-TEMPLATES
% ==========================================

\cs_new_protected:Npn \__ltxtalk_apply_class_templates:
  {
    \EditInstance { header } { std }
      {
        background-color = ltxtalk-primary,
        color = white,
        font = \bfseries,
        height = \headheight + \headsep,
        print-frame-title = true
      }
    \EditInstance { footer } { std }
      {
        background-color = ltxtalk-bg,
        color = ltxtalk-secondary,
        element-order = { author, institute, date, framenumber },
        font = \tiny
      }
    \EditInstance { frametitle } { header }
      {
        color = white,
        font = \Large \bfseries,
        before-vspace = 0pt,
        after-vspace = 0pt
      }
    \EditInstance { titlepage-element } { title }
      { color = ltxtalk-primary, font = \Large \bfseries }
    \EditInstance { titlepage-element } { subtitle }
      { color = ltxtalk-secondary }
    \str_case:Vn \g__ltxtalk_active_theme_tl
      {
        { minimal }
          {
            \EditInstance { header } { std }
              { background-color =, color = ltxtalk-primary }
            \EditInstance { footer } { std }
              { background-color =, element-order = { framenumber } }
            \EditInstance { frametitle } { header }
              { color = ltxtalk-primary, font = \Large \bfseries }
          }
        { academic }
          {
            \EditInstance { header } { std }
              { background-color = ltxtalk-primary, color = white }
            \EditInstance { footer } { std }
              { element-order = { author, institute, framenumber } }
          }
        { corporate }
          {
            \EditInstance { header } { std }
              { background-color = ltxtalk-primary, color = white }
            \EditInstance { footer } { std }
              {
                background-color = ltxtalk-primary,
                color = white,
                element-order = { institute, date, framenumber }
              }
          }
      }
  }

% ==========================================
% INITIALISIERUNG
% ==========================================

\AtBeginDocument{
  % Prüfen ob Themes geladen aber nicht aktiviert wurden
  \str_if_eq:VnT \g__ltxtalk_active_theme_tl { default }
    {
      \clist_map_inline:nn { modern, academic, corporate, minimal }
        {
          \@ifpackageloaded{ltxtalk-theme-#1}
            {
              \msg_warning:nnn { ltxtalk } { theme-not-activated } { #1 }
            }{}
        }
    }
}

% ==========================================
% GENERIC PAGE-STYLE ADAPTER
% ==========================================
% \ldeen{Als Schnittstelle zu \pkg*{ltx-talk} dient der Seitenstil.
% }{
% The page style constitutes the interface to \pkg*{ltx-talk}.}
\seq_new:N \g__ltxtalk_header_slots_seq
\seq_new:N \g__ltxtalk_left_slots_seq
\seq_new:N \g__ltxtalk_content_slots_seq
\seq_new:N \g__ltxtalk_right_slots_seq
\seq_new:N \g__ltxtalk_footer_slots_seq
\prop_new:N \g__ltxtalk_slot_contents_prop

\coffin_new:N \l__ltxtalk_area_coffin
\coffin_new:N \l__ltxtalk_slot_coffin
\dim_new:N \l__ltxtalk_current_area_width_dim
\dim_new:N \l__ltxtalk_current_area_height_dim
\dim_new:N \l__ltxtalk_current_slot_width_dim
\dim_new:N \l__ltxtalk_current_slot_height_dim
\dim_new:N \l__ltxtalk_current_slot_x_dim
\dim_new:N \l__ltxtalk_current_slot_y_dim
\dim_new:N \l__ltxtalk_grid_star_width_dim
\dim_new:N \l__ltxtalk_grid_column_width_dim
\tl_new:N \l__ltxtalk_current_area_tl
\int_new:N \l__ltxtalk_grid_star_count_int
\dim_new:N \g__ltxtalk_base_left_margin_dim
\dim_new:N \g__ltxtalk_base_right_margin_dim
\dim_new:N \g__ltxtalk_current_left_margin_dim
\dim_new:N \g__ltxtalk_current_right_margin_dim
\dim_gset:Nn \g__ltxtalk_base_left_margin_dim {\Gm@lmargin}
\dim_gset:Nn \g__ltxtalk_base_right_margin_dim {\Gm@rmargin}

\NewExpandableDocumentCommand \LTXTalkBaseLeftMargin {}
  {\dim_use:N \g__ltxtalk_base_left_margin_dim}
\NewExpandableDocumentCommand \LTXTalkBaseRightMargin {}
  {\dim_use:N \g__ltxtalk_base_right_margin_dim}
\NewExpandableDocumentCommand \LTXTalkCurrentLeftMargin {}
  {\dim_use:N \g__ltxtalk_current_left_margin_dim}
\NewExpandableDocumentCommand \LTXTalkCurrentRightMargin {}
  {\dim_use:N \g__ltxtalk_current_right_margin_dim}

\NewTemplateType { lttheme-slot } { 1 }
\DeclareTemplateInterface { lttheme-slot } { standard } { 1 }
  {
    font        : tokenlist = \normalfont,
    color       : tokenlist = ltxtalk-text,
    background  : tokenlist = none,
    align       : choice { left, center, right } = center,
    frame       : choice { none, box, top, bottom, left, right } = none,
    padding     : length = 0pt,
    rule-color  : tokenlist = ltxtalk-primary,
    rule-thick  : length = 0.4pt,
    rotate      : tokenlist = 0,
    opacity     : tokenlist = 1,
    shadow      : boolean = false,
    uppercase   : boolean = false,
    lowercase   : boolean = false
  }
\DeclareTemplateCode { lttheme-slot } { standard } { 1 }
  {
    font       = \l__ltxtalk_slot_font_tl,
    color      = \l__ltxtalk_slot_color_tl,
    background = \l__ltxtalk_slot_bg_tl,
    align      =
      {
        left   = \tl_set:Nn \l__ltxtalk_slot_align_tl { left },
        center = \tl_set:Nn \l__ltxtalk_slot_align_tl { center },
        right  = \tl_set:Nn \l__ltxtalk_slot_align_tl { right }
      },
    frame      =
      {
        none   = \tl_set:Nn \l__ltxtalk_slot_frame_tl { none },
        box    = \tl_set:Nn \l__ltxtalk_slot_frame_tl { box },
        top    = \tl_set:Nn \l__ltxtalk_slot_frame_tl { top },
        bottom = \tl_set:Nn \l__ltxtalk_slot_frame_tl { bottom },
        left   = \tl_set:Nn \l__ltxtalk_slot_frame_tl { left },
        right  = \tl_set:Nn \l__ltxtalk_slot_frame_tl { right }
      },
    padding    = \l__ltxtalk_slot_padding_dim,
    rule-color = \l__ltxtalk_slot_rule_tl,
    rule-thick = \l__ltxtalk_slot_rule_thick_dim,
    rotate = \l__ltxtalk_slot_rotate_tl,
    opacity = \l__ltxtalk_slot_opacity_tl,
    shadow = \l__ltxtalk_slot_shadow_bool,
    uppercase = \l__ltxtalk_slot_uppercase_bool,
    lowercase = \l__ltxtalk_slot_lowercase_bool
  }
  { \__ltxtalk_typeset_slot:n {#1} }

\cs_new_protected:Npn \__ltxtalk_typeset_slot:n #1
  {
    \group_begin:
      \setlength\fboxsep{\l__ltxtalk_slot_padding_dim}
      \setlength\fboxrule{\l__ltxtalk_slot_rule_thick_dim}
      \color{\l__ltxtalk_slot_color_tl}
      \l__ltxtalk_slot_font_tl
      \str_case:Vn \l__ltxtalk_slot_align_tl
        {
          { left }   { \raggedright }
          { center } { \centering }
          { right }  { \raggedleft }
        }
      \str_if_eq:VnF \l__ltxtalk_slot_frame_tl { none }
        { \color{\l__ltxtalk_slot_rule_tl} }
      \str_if_eq:VnTF \l__ltxtalk_slot_frame_tl { box }
        {
          \fcolorbox{\l__ltxtalk_slot_rule_tl}
            {\str_if_eq:VnTF \l__ltxtalk_slot_bg_tl {none}{white}{\l__ltxtalk_slot_bg_tl}}
            {\color{\l__ltxtalk_slot_color_tl}#1}
        }
        {
          \str_if_eq:VnF \l__ltxtalk_slot_bg_tl { none }
            { \colorbox{\l__ltxtalk_slot_bg_tl}{\color{\l__ltxtalk_slot_color_tl}#1} }
          \str_if_eq:VnT \l__ltxtalk_slot_bg_tl { none } {#1}
        }
      \par
    \group_end:
  }

\NewTemplateType { lttheme-area } { 1 }
\DeclareTemplateInterface { lttheme-area } { standard } { 1 }
  {
    background     : tokenlist = none,
    frame          : choice { none, box, top, bottom, left, right, topbottom } = none,
    rule-color     : tokenlist = ltxtalk-primary,
    rule-thickness : length = 0.4pt,
    padding        : length = 0pt,
    artifact-type  : tokenlist = layout,
    background-opacity : tokenlist = 1,
    shadow         : boolean = false,
    shadow-color   : tokenlist = black!30,
    shadow-offset  : length = 3pt,
    margin-top     : length = 0pt,
    margin-bottom  : length = 0pt
  }
\DeclareTemplateCode { lttheme-area } { standard } { 1 }
  {
    background = \l__ltxtalk_area_background_tl,
    frame =
      {
        none      = \tl_set:Nn \l__ltxtalk_area_frame_tl { none },
        box       = \tl_set:Nn \l__ltxtalk_area_frame_tl { box },
        top       = \tl_set:Nn \l__ltxtalk_area_frame_tl { top },
        bottom    = \tl_set:Nn \l__ltxtalk_area_frame_tl { bottom },
        left      = \tl_set:Nn \l__ltxtalk_area_frame_tl { left },
        right     = \tl_set:Nn \l__ltxtalk_area_frame_tl { right },
        topbottom = \tl_set:Nn \l__ltxtalk_area_frame_tl { topbottom }
      },
    rule-color = \l__ltxtalk_area_rule_color_tl,
    rule-thickness = \l__ltxtalk_area_rule_thickness_dim,
    padding = \l__ltxtalk_area_padding_dim,
    artifact-type = \l__ltxtalk_area_artifact_tl,
    background-opacity = \l__ltxtalk_area_bg_opacity_tl,
    shadow = \l__ltxtalk_area_shadow_bool,
    shadow-color = \l__ltxtalk_area_shadow_color_tl,
    shadow-offset = \l__ltxtalk_area_shadow_offset_dim,
    margin-top = \l__ltxtalk_area_margin_top_dim,
    margin-bottom = \l__ltxtalk_area_margin_bottom_dim
  }
  { \__ltxtalk_typeset_area_background:n {#1} }

\clist_map_inline:nn { header, left, content, right, footer }
  { \DeclareInstance { lttheme-area } {#1} { standard } { } }

\cs_new_protected:Npn \__ltxtalk_typeset_area_background:n #1
  {
    \group_begin:
      \setlength\fboxsep{0pt}
      \clist_if_in:nVT {top,topbottom,box} \l__ltxtalk_area_frame_tl
        {
          \color{\l__ltxtalk_area_rule_color_tl}
          \rule{\l__ltxtalk_current_area_width_dim}
            {\l__ltxtalk_area_rule_thickness_dim}\par\nointerlineskip
        }
      \str_if_eq:VnTF \l__ltxtalk_area_background_tl { none }
        {\parbox[c][\l__ltxtalk_current_area_height_dim][c]
           {\l__ltxtalk_current_area_width_dim}
           {\__ltxtalk_typeset_area_sides:n {#1}}}
        {\colorbox{\l__ltxtalk_area_background_tl}
           {\parbox[c][\l__ltxtalk_current_area_height_dim][c]
             {\l__ltxtalk_current_area_width_dim}
             {\__ltxtalk_typeset_area_sides:n {#1}}}}
      \clist_if_in:nVT {bottom,topbottom,box} \l__ltxtalk_area_frame_tl
        {
          \par\nointerlineskip\color{\l__ltxtalk_area_rule_color_tl}
          \rule{\l__ltxtalk_current_area_width_dim}
            {\l__ltxtalk_area_rule_thickness_dim}
        }
    \group_end:
  }
\cs_new_protected:Npn \__ltxtalk_typeset_area_sides:n #1
  {
    \str_if_eq:VnT \l__ltxtalk_area_frame_tl {left}
      {
        \color{\l__ltxtalk_area_rule_color_tl}
        \rule{\l__ltxtalk_area_rule_thickness_dim}
          {\l__ltxtalk_current_area_height_dim}\hspace{\l__ltxtalk_area_padding_dim}
      }
    #1
    \str_if_eq:VnT \l__ltxtalk_area_frame_tl {right}
      {
        \hfill\color{\l__ltxtalk_area_rule_color_tl}
        \rule{\l__ltxtalk_area_rule_thickness_dim}
          {\l__ltxtalk_current_area_height_dim}
      }
  }
% \ldeen{Über Sockets können die einzelnen Bereichebehandlungen gezielt ersetzt werden.
% }{
% The different area handler can be specificaly replcaced by means of sockets.}
\NewSocket { lttheme/page/background } { 0 }
\NewSocket { lttheme/area/header } { 0 }
\NewSocket { lttheme/area/left } { 0 }
\NewSocket { lttheme/area/right } { 0 }
\NewSocket { lttheme/area/footer } { 0 }
\NewSocket { lttheme/page/foreground } { 0 }
\NewSocketPlug { lttheme/page/background } { standard } { }
\NewSocketPlug { lttheme/page/background } { magpie-title }
  {
    \int_compare:nNnT {\@framenumber} = {1}
      {
        \tag_mc_begin:n {artifact}
        \raisebox{-2.05cm}[0pt][0pt]
          {\color{ltxtalk-primary}\rule{\paperwidth}{1.55cm}}
        \tag_mc_end:
      }
  }
\NewSocketPlug { lttheme/area/header } { standard }
  { \__ltxtalk_render_area:nnn {header}{\paperwidth}{\l__ltxtalk_header_height_dim} }
\NewSocketPlug { lttheme/area/left } { standard }
  { \__ltxtalk_render_area:nnn {left}{\l__ltxtalk_left_area_width_dim}{\textheight} }
\NewSocketPlug { lttheme/area/left } { hawk }
  { \__ltxtalk_render_area:nnn {left}{\l__ltxtalk_left_area_width_dim}{\paperheight} }
\NewSocketPlug { lttheme/area/right } { standard }
  { \__ltxtalk_render_area:nnn {right}{\l__ltxtalk_right_area_width_dim}{\textheight} }
\NewSocketPlug { lttheme/area/footer } { standard }
  { \__ltxtalk_render_area:nnn {footer}{\paperwidth}{\l__ltxtalk_footer_height_dim} }
\NewSocketPlug { lttheme/page/foreground } { standard } { }
\NewSocketPlug { lttheme/page/foreground } { magpie-title }
  {
    \int_compare:nNnT {\@framenumber} = {1}
      {
        \tag_mc_begin:n {artifact}
        \raisebox{-1.42cm}[0pt][0pt]
          {\parbox[c][1.25cm][c]{\paperwidth}
            {\centering\color{white}\sffamily
             \Large\bfseries\@title\par
             \small\mdseries\@subtitle}}
        \tag_mc_end:
      }
  }
\NewSocketPlug { lttheme/page/foreground } { sparrow-bars }
  {
    \tag_mc_begin:n {artifact}
    \raisebox{\Gm@tmargin}[0pt][0pt]
      {\vbox to 0pt
        {
          \color{ltxtalk-primary!35}\rlap{\rule{\paperwidth}{0.35pt}}
          \vskip 0.42cm
          \color{ltxtalk-primary!20}\rlap{\rule{\paperwidth}{0.35pt}}
          \vskip 0.30cm
          \color{ltxtalk-primary!45}\rlap{\rule{\paperwidth}{0.7pt}}
          \vss
        }}
    \tag_mc_end:
  }
\clist_map_inline:nn
  { page/background, area/header, area/left, area/right, area/footer, page/foreground }
  { \AssignSocketPlug { lttheme/#1 } { standard } }

\cs_new_protected:Npn \__ltxtalk_clear_theme_state:
  {
    \AssignSocketPlug { lttheme/page/background } { standard }
    \AssignSocketPlug { lttheme/area/header } { standard }
    \AssignSocketPlug { lttheme/area/left } { standard }
    \AssignSocketPlug { lttheme/area/right } { standard }
    \AssignSocketPlug { lttheme/area/footer } { standard }
    \AssignSocketPlug { lttheme/page/foreground } { standard }
    \prop_gclear:N \g__ltxtalk_slot_types_prop
    \prop_gclear:N \g__ltxtalk_slot_positions_prop
    \prop_gclear:N \g__ltxtalk_slot_styles_prop
    \prop_gclear:N \g__ltxtalk_slot_areas_prop
    \prop_gclear:N \g__ltxtalk_slot_accessibility_prop
    \prop_gclear:N \g__ltxtalk_area_grid_columns_prop
    \prop_gclear:N \g__ltxtalk_area_grid_rows_prop
    \prop_gclear:N \g__ltxtalk_slot_grid_positions_prop
    \prop_gclear:N \g__ltxtalk_slot_contents_prop
    \prop_gclear:N \g__ltxtalk_area_decorations_prop
    \seq_gclear:N \g__ltxtalk_header_slots_seq
    \seq_gclear:N \g__ltxtalk_left_slots_seq
    \seq_gclear:N \g__ltxtalk_content_slots_seq
    \seq_gclear:N \g__ltxtalk_right_slots_seq
    \seq_gclear:N \g__ltxtalk_footer_slots_seq
  }

\cs_new_protected:Npn \__ltxtalk_update_text_area:
  {
    \dim_gset_eq:NN \g__ltxtalk_current_left_margin_dim
      \g__ltxtalk_base_left_margin_dim
    \dim_gset_eq:NN \g__ltxtalk_current_right_margin_dim
      \g__ltxtalk_base_right_margin_dim
    \bool_if:NT \l__ltxtalk_show_left_margin_bool
      {
        \dim_gadd:Nn \g__ltxtalk_current_left_margin_dim
          {\l__ltxtalk_content_inset_left_dim}
      }
    \bool_if:NT \l__ltxtalk_show_right_margin_bool
      {
        \dim_gadd:Nn \g__ltxtalk_current_right_margin_dim
          {\l__ltxtalk_content_inset_right_dim}
      }
    \cs_gset:Npx \Gm@lmargin
      {\dim_use:N \g__ltxtalk_current_left_margin_dim}
    \cs_gset:Npx \Gm@rmargin
      {\dim_use:N \g__ltxtalk_current_right_margin_dim}
    \dim_gset:Nn \textwidth
      {
        \paperwidth-\g__ltxtalk_current_left_margin_dim
        -\g__ltxtalk_current_right_margin_dim
      }
    \dim_gset:Nn \oddsidemargin
      {\g__ltxtalk_current_left_margin_dim-1in}
    \dim_gset_eq:NN \evensidemargin \oddsidemargin
    \dim_gset_eq:NN \columnwidth \textwidth
    \dim_gset_eq:NN \linewidth \textwidth
    \dim_gset_eq:NN \hsize \textwidth
  }

\RenewDocumentCommand \defineltxtalkslot { m m m m } {
  \prop_gput:Nnn \g__ltxtalk_slot_areas_prop {#1} {#2}
  \prop_gput:Nnn \g__ltxtalk_slot_types_prop {#1} {#3}
  \prop_gput:Nnn \g__ltxtalk_slot_accessibility_prop {#1} {#4}
  \seq_if_exist:cF { g__ltxtalk_#2_slots_seq }
    { \seq_new:c { g__ltxtalk_#2_slots_seq } }
  \seq_gput_right:cn { g__ltxtalk_#2_slots_seq } {#1}
  \IfInstanceExistsTF { lttheme-slot } {#1}
    { \EditInstance { lttheme-slot } {#1} { } }
    { \DeclareInstance { lttheme-slot } {#1} { standard } { } }
}

\NewDocumentCommand \setltxtalkslotcontent { m m } {
  \prop_gput:Nnn \g__ltxtalk_slot_contents_prop {#1} {#2} 
}

\RenewDocumentCommand \styleltxtalkslot { m m } {
  \prop_gput:Nnn \g__ltxtalk_slot_styles_prop {#1} {#2}
  \EditInstance { lttheme-slot } {#1} {#2}
}

\RenewDocumentCommand \decorateltxtalkarea { m m } {
  \prop_gput:Nnn \g__ltxtalk_area_decorations_prop {#1} {#2}
  \EditInstance { lttheme-area } {#1} {#2}
}

\RenewDocumentCommand \useltxtalktheme { m } {
  \str_if_eq:nnTF {#1} {default}
    {
      \__ltxtalk_clear_theme_state:
      \tl_gset:Nn \g__ltxtalk_active_theme_tl {default}
      \__ltxtalk_update_text_area:
      \pagestyle{talk}
    }
    {
      \prop_get:NnNTF \g__ltxtalk_themes_prop {#1} \l_tmpa_tl
        {
          \__ltxtalk_clear_theme_state:
          \tl_gset:Nn \g__ltxtalk_active_theme_tl {#1}
          \l_tmpa_tl
          \__ltxtalk_apply_class_templates:
        }
        { \msg_error:nnn {ltxtalk}{unknown-theme}{#1} }
    }
}

\NewDocumentCommand \addltxtalktheme { m } {
  \prop_get:NnNTF \g__ltxtalk_themes_prop {#1} \l_tmpa_tl
    { \l_tmpa_tl \__ltxtalk_apply_class_templates: }
    { \msg_error:nnn {ltxtalk}{unknown-theme}{#1} }
}

\cs_new_protected:Npn \__ltxtalk_slot_content:n #1
  {
    \prop_get:NnN \g__ltxtalk_slot_types_prop {#1} \l__ltxtalk_slot_type_tl
    \str_case:VnF \l__ltxtalk_slot_type_tl
      {
        {title}         {\LTXTalkMetadata{short-title}}
        {subtitle}      {\LTXTalkMetadata{short-subtitle}}
        {author}        {\LTXTalkMetadata{short-author}}
        {date}          {\LTXTalkMetadata{date}}
        {institute}     {\LTXTalkMetadata{short-institute}}
        {section}       {\LTXTalkMetadata{section}}
        {subsection}    {\LTXTalkMetadata{subsection}}
        {frametitle}
          {\__talk_frame_title:n {\tl_use:N \g__talk_frame_title_tl}}
        {framesubtitle} {\tl_use:N \g__talk_frame_subtitle_tl}
        {pagenumber}    {\LTXTalkMetadata{framenumber}}
      }
      { \prop_item:Nn \g__ltxtalk_slot_contents_prop {#1} }
  }

\cs_new_protected:Npn \__ltxtalk_render_slot:n #1
  {
    \prop_get:NnNTF \g__ltxtalk_slot_grid_positions_prop {#1} \l_tmpa_tl
      {\exp_after:wN \__ltxtalk_render_grid_slot_aux:w
        \l_tmpa_tl \q_stop {#1}}
      {\prop_get:NnNT \g__ltxtalk_slot_positions_prop {#1} \l_tmpa_tl
        {\exp_after:wN \__ltxtalk_render_slot_aux:w
          \l_tmpa_tl \q_stop {#1}}}
  }
\cs_new_protected:Npn \__ltxtalk_render_grid_slot_aux:w
  #1#2#3#4 \q_stop #5
  {\__ltxtalk_render_grid_slot:nnnnn {#5}{#1}{#2}{#3}{#4}}
\cs_new_protected:Npn \__ltxtalk_render_grid_slot:nnnnn #1#2#3#4#5
  {
    \prop_get:NVN \g__ltxtalk_area_grid_columns_prop
      \l__ltxtalk_current_area_tl \l_tmpa_tl
    \prop_get:NVN \g__ltxtalk_area_grid_rows_prop
      \l__ltxtalk_current_area_tl \l_tmpb_tl
    \clist_set:NV \l_tmpa_clist \l_tmpa_tl
    \dim_zero:N \l_tmpa_dim
    \int_zero:N \l__ltxtalk_grid_star_count_int
    \clist_map_inline:Nn \l_tmpa_clist
      {\str_if_eq:nnTF {##1}{*}
        {\int_incr:N \l__ltxtalk_grid_star_count_int}
        {\dim_add:Nn \l_tmpa_dim {##1}}}
    \int_compare:nNnTF {\l__ltxtalk_grid_star_count_int} > {0}
      {
        \dim_set_eq:NN \l__ltxtalk_grid_star_width_dim
          \l__ltxtalk_current_area_width_dim
        \dim_sub:Nn \l__ltxtalk_grid_star_width_dim {\l_tmpa_dim}
        \dim_set:Nn \l__ltxtalk_grid_star_width_dim
          {\l__ltxtalk_grid_star_width_dim/\l__ltxtalk_grid_star_count_int}
      }
      {\dim_zero:N \l__ltxtalk_grid_star_width_dim}
    \dim_zero:N \l__ltxtalk_current_slot_x_dim
    \dim_zero:N \l__ltxtalk_current_slot_width_dim
    \int_zero:N \l_tmpa_int
    \clist_map_inline:Nn \l_tmpa_clist
      {
        \int_incr:N \l_tmpa_int
        \str_if_eq:nnTF {##1}{*}
          {\dim_set_eq:NN \l__ltxtalk_grid_column_width_dim
             \l__ltxtalk_grid_star_width_dim}
          {\dim_set:Nn \l__ltxtalk_grid_column_width_dim {##1}}
        \int_compare:nNnT {\l_tmpa_int} < {#2}
          {\dim_add:Nn \l__ltxtalk_current_slot_x_dim
            {\l__ltxtalk_grid_column_width_dim}}
        \int_compare:nT {\l_tmpa_int >= #2}
          {\int_compare:nT {\l_tmpa_int < #2+#4}
            {\dim_add:Nn \l__ltxtalk_current_slot_width_dim
              {\l__ltxtalk_grid_column_width_dim}}}
      }
    \dim_set:Nn \l_tmpa_dim
      {\l__ltxtalk_current_area_height_dim/\l_tmpb_tl}
    \dim_set:Nn \l__ltxtalk_current_slot_y_dim
      {\int_eval:n {#3-1}\l_tmpa_dim}
    \dim_set:Nn \l__ltxtalk_current_slot_height_dim
      {#5\l_tmpa_dim}
    \__ltxtalk_render_slot:nnnnn {#1}
      {\l__ltxtalk_current_slot_x_dim}{\l__ltxtalk_current_slot_y_dim}
      {\l__ltxtalk_current_slot_width_dim}{\l__ltxtalk_current_slot_height_dim}
  }
\cs_new_protected:Npn \__ltxtalk_render_slot_aux:w #1#2#3#4 \q_stop #5
  { \__ltxtalk_render_slot:nnnnn {#5}{#1}{#2}{#3}{#4} }
\cs_new_protected:Npn \__ltxtalk_render_slot:nnnnn #1#2#3#4#5
  {
    \dim_set:Nn \l__ltxtalk_current_slot_width_dim {#4}
    \dim_set:Nn \l__ltxtalk_current_slot_height_dim {#5}
    \dim_set:Nn \l__ltxtalk_current_slot_x_dim
      {
        \str_case:nnF {#2}
          {
            {left}{0pt}
            {center}{(\l__ltxtalk_current_area_width_dim-#4)/2}
            {right}{\l__ltxtalk_current_area_width_dim-#4}
          }{#2}
      }
    \dim_set:Nn \l__ltxtalk_current_slot_y_dim
      {
        \str_case:nnF {#3}
          {
            {top}{0pt}
            {center}{(\l__ltxtalk_current_area_height_dim-#5)/2}
            {bottom}{\l__ltxtalk_current_area_height_dim-#5}
          }{#3}
      }
    \hcoffin_set:Nn \l__ltxtalk_slot_coffin
      {
        \parbox[c][#5][c]{#4}
          {\UseInstance {lttheme-slot}{#1}{\__ltxtalk_slot_content:n {#1}}}
      }
    \coffin_join:NnnNnnnn \l__ltxtalk_area_coffin {l}{t}
      \l__ltxtalk_slot_coffin {l}{t}
      {\l__ltxtalk_current_slot_x_dim}{-\l__ltxtalk_current_slot_y_dim}
  }

\cs_new_protected:Npn \__ltxtalk_render_area:nnn #1#2#3
  {
    \tl_set:Nn \l__ltxtalk_current_area_tl {#1}
    \dim_set:Nn \l__ltxtalk_current_area_width_dim {#2}
    \dim_set:Nn \l__ltxtalk_current_area_height_dim {#3}
    \vcoffin_set:Nnn \l__ltxtalk_area_coffin {#2}
      {\UseInstance {lttheme-area}{#1}{\rule{0pt}{#3}}}
    \seq_map_inline:cn {g__ltxtalk_#1_slots_seq}
      {\__ltxtalk_render_slot:n {##1}}
    \coffin_typeset:Nnnnn \l__ltxtalk_area_coffin {l}{t}{0pt}{0pt}
  }
% \ldeen{
%   Inhalts-Slots gehören zum Textstrom des Frames und nicht zum Seitenlayout.
%   Dieser schmale Adapter behält das Klassenkommando bei und fügt alle passenden 
%   Slots in der Reihenfolge ihrer Deklaration hinzu. Folglich kann ein Theme einen 
%   Titel im Hauptblock, in der Kopfzeile oder in beiden Bereichen platzieren, ohne 
%   dass die Themenengine das Theme kennt.
% }{
%   Content slots belong to the frame's text stream, not to the page layout.
%   This narrow adapter preserves the class command and adds all matching slots
%   in declaration order.  Consequently a theme may place a title in the main
%   block, in the header, or in both areas without the theme engine knowing the theme.
% }
\cs_new_protected:Npn \__ltxtalk_render_content_slots:n #1
  {
    \seq_map_inline:Nn \g__ltxtalk_content_slots_seq
      {
        \prop_get:NnNT \g__ltxtalk_slot_types_prop {##1} \l_tmpa_tl
          {
            \str_if_eq:VnT \l_tmpa_tl {#1}
              {
                \prop_get:NnNT \g__ltxtalk_slot_positions_prop {##1} \l_tmpb_tl
                  {
                    \exp_after:wN \__ltxtalk_render_content_slot_aux:w
                      \l_tmpb_tl \q_stop {##1}
                  }
              }
          }
      }
  }
% \ldeen{Kein führendes \cs{par}\cs{noindent} vor dem Slot-Inhalt: Für den
%   Slot-Typ \code{frametitle} ruft \cs{__ltxtalk_slot_content:n} die
%   klasseneigene Funktion \cs{__talk_frame_title:n} auf, die bei aktivem
%   Tagging den Inhalt explizit taggt und dafür \cs{tagpdfparaOff} lokal
%   setzt. Ein vorangehender \cs{par} löst jedoch bereits den
%   automatischen Absatz-Tagging-Hook von tagpdf für einen (leeren) Absatz
%   aus, bevor \cs{tagpdfparaOff} greift; dieser Hook bleibt dann
%   unausgeglichen offen und tagpdf bricht den Lauf mit \emph{"the number
%   of automatic begin and end text-unit para hooks differ"} fatal ab.
%   Das schließende \cs{par} nach dem Inhalt bleibt nötig, damit der Slot
%   sauber als eigener Absatz endet, bevor nachfolgender Frame-Inhalt
%   beginnt.}{
%   No leading \cs{par}\cs{noindent} before the slot content: for the
%   \code{frametitle} slot type, \cs{__ltxtalk_slot_content:n} calls the
%   class's own \cs{__talk_frame_title:n}, which explicitly tags the
%   content when tagging is active and locally sets \cs{tagpdfparaOff}
%   for that. A preceding \cs{par}, however, already fires tagpdf's
%   automatic per-paragraph tagging hook for an (empty) paragraph before
%   \cs{tagpdfparaOff} takes effect; that hook then stays unbalanced and
%   tagpdf aborts the run fatally with \emph{"the number of automatic
%   begin and end text-unit para hooks differ"}. The trailing \cs{par}
%   after the content remains necessary so the slot cleanly ends its own
%   paragraph before subsequent frame content begins.}
\cs_new_protected:Npn \__ltxtalk_render_content_slot_aux:w
  #1#2#3#4 \q_stop #5
  {
    \UseInstance {lttheme-slot}{#5}{\__ltxtalk_slot_content:n {#5}}
    \par
  }
\RenewDocumentCommand \frametitle { D <> { all } O {#3} m }
  {
    \__talk_if_overlay:nT {#1}
      {
        \tl_gset:Nn \g__talk_frame_title_tl {#3}
        \__ltxtalk_render_content_slots:n {frametitle}
      }
  }
\RenewDocumentCommand \framesubtitle { D <> { all } O {#3} m }
  {
    \__talk_if_overlay:nT {#1}
      {
        \tl_gset:Nn \g__talk_frame_subtitle_tl {#3}
        \__ltxtalk_render_content_slots:n {framesubtitle}
      }
  }

\cs_new_nopar:Npn \ps@lttheme
  {
    \cs_set_nopar:Npn \@oddhead
      {
        \hbox_to_wd:nn {0pt}
          {
            \skip_horizontal:n {-\Gm@lmargin}
            \UseSocket{lttheme/page/background}
            \hss
          }
        \hbox_to_wd:nn {0pt}
          {
            \skip_horizontal:n {-\Gm@lmargin}
            \bool_if:NT \l__ltxtalk_show_left_margin_bool
              {
                \str_if_eq:VnTF \g__ltxtalk_active_theme_tl {hawk}
                  { \raisebox{\Gm@tmargin}[0pt][0pt]{\UseSocket{lttheme/area/left}} }
                  {
                    \raisebox
                      {-\dim_eval:n{\l__ltxtalk_header_height_dim-\Gm@tmargin}}
                      [0pt][0pt]{\UseSocket{lttheme/area/left}}
                  }
              }
            \hss
          }
        \hbox_to_wd:nn {0pt}
          {
            \bool_if:NT \l__ltxtalk_show_right_margin_bool
              {
                \skip_horizontal:n
                  {
                    \dim_eval:n
                      {-\Gm@lmargin+\paperwidth-\l__ltxtalk_right_area_width_dim}
                  }
                \raisebox
                  {-\dim_eval:n{\l__ltxtalk_header_height_dim-\Gm@tmargin}}
                  [0pt][0pt]
                  {\UseSocket{lttheme/area/right}}
              }
            \hss
          }
        \hbox_to_wd:nn {0pt}
          {
            \bool_if:NT \l__ltxtalk_show_header_bool
              {
                \str_if_eq:VnTF \g__ltxtalk_active_theme_tl {magpie}
                  {
                    \int_compare:nNnF {\@framenumber} = {1}
                      {
                        \skip_horizontal:n {-\Gm@lmargin}
                        \raisebox{\Gm@tmargin}[0pt][0pt]
                          {\UseSocket{lttheme/area/header}}
                      }
                  }
                  {
                    \skip_horizontal:n {-\Gm@lmargin}
                    \raisebox{\Gm@tmargin}[0pt][0pt]
                      {\UseSocket{lttheme/area/header}}
                  }
              }
            \hss
          }
        \hbox_to_wd:nn {0pt}
          {
            \skip_horizontal:n {-\Gm@lmargin}
            \UseSocket{lttheme/page/foreground}
            \hss
          }
        \hfil
      }
    \cs_set_nopar:Npn \@oddfoot
      {
        \hbox_to_wd:nn {0pt}
          {
            \bool_if:NT \l__ltxtalk_show_footer_bool
              {
                \skip_horizontal:n {-\Gm@lmargin}
                \raisebox
                  {
                    \dim_eval:n
                      {
                        \l__ltxtalk_footer_height_dim
                        -\Gm@bmargin+\footskip
                      }
                  }
                  [0pt][0pt]
                  {\UseSocket{lttheme/area/footer}}
              }
            \hss
          }
        \hfil
      }
    \cs_set_eq:NN \@evenhead \@oddhead
    \cs_set_eq:NN \@evenfoot \@oddfoot
  }

\cs_set_protected:Npn \__ltxtalk_apply_class_templates:
  {
    \__ltxtalk_update_text_area:
    \pagestyle{lttheme}
  }
\AddToHook { begindocument / end }
  {
    \str_if_eq:VnF \g__ltxtalk_active_theme_tl {default}
      {\__ltxtalk_update_text_area:}
  }

%</pkg>
%<*minimal>
% \subsection{Minimal theme}
\NeedsTeXFormat{LaTeX2e}
\ProvidesExplPackage{ltxtalk-theme-minimal}
{2026/08/05}
{v0.4.1}
{Minimal theme for ltx-talk}

\RequirePackage{ltxtalk-theme}

\cs_new_protected:Npn \__ltxtalk_minimal_theme_setup:
  {
    % === LAYOUT ===
    \ltxtalksetup{
      header = true,
      header-height = 1.2cm,
      footer = true,
      footer-height = 0.6cm,
    }
    
    % === TITEL-PLATZIERUNG ===
    \setltxtalktitleplacement{
      frametitle = content,
      section = none,
    }
    
    % === SLOTS ===
    \defineltxtalkslot{min-title}{header}{title}{
      role = note
    }
    \defineltxtalkslot{min-section}{header}{section}{
      role = note
    }
    \defineltxtalkslot{min-page}{header}{pagenumber}{
      role = navigation
    }
    \defineltxtalkslot{min-ftitle}{content}{frametitle}{
      role = heading, heading-level = 2
    }
    \defineltxtalkslot{min-fsubtitle}{content}{framesubtitle}{
      role = heading, heading-level = 3
    }
    \defineltxtalkslot{min-author}{footer}{author}{
      role = note
    }
    \defineltxtalkslot{min-footer-page}{footer}{pagenumber}{
      role = navigation
    }
    
    % === POSITIONIERUNG ===
    \positionltxtalkslot{min-title}{0.18cm}{0.3cm}{0.4\paperwidth}{0.6cm}
    \positionltxtalkslot{min-section}{center}{0.3cm}{0.3\paperwidth}{0.6cm}
    \positionltxtalkslot{min-page}{right}{0.3cm}{0.1\paperwidth}{0.6cm}
    \positionltxtalkslot{min-ftitle}{left}{1cm}{0.9\paperwidth}{2cm}
    \positionltxtalkslot{min-fsubtitle}{left}{3cm}{0.9\paperwidth}{1cm}
    \positionltxtalkslot{min-author}{0.18cm}{bottom}{0.4\paperwidth}{0.4cm}
    \positionltxtalkslot{min-footer-page}
      {\dim_eval:n{0.9\paperwidth-0.18cm}}{bottom}{0.1\paperwidth}{0.4cm}
    
    % === FORMATIERUNG ===
    \styleltxtalkslot{min-title}{
      font = \small\sffamily,
      color = ltxtalk-secondary,
      align = left,
    }
    \styleltxtalkslot{min-section}{
      font = \small\sffamily,
      color = ltxtalk-secondary!70,
    }
    \styleltxtalkslot{min-page}{
      font = \tiny\sffamily,
      color = ltxtalk-secondary!50,
      align = right,
    }
    \styleltxtalkslot{min-ftitle}{
      font = \LARGE\bfseries,
      color = ltxtalk-primary,
      align = left,
    }
    \styleltxtalkslot{min-fsubtitle}{
      font = \large\itshape,
      color = ltxtalk-secondary,
      align = left,
    }
    \styleltxtalkslot{min-author}{
      font = \tiny\itshape,
      color = ltxtalk-secondary!60,
      align = left,
    }
    \styleltxtalkslot{min-footer-page}{
      font = \tiny\sffamily,
      color = ltxtalk-secondary!50,
      align = right,
    }
    
    % === BEREICHSDEKORATIONEN ===
    \decorateltxtalkarea{header}{
      frame = bottom,
      rule-color = ltxtalk-primary,
      rule-thickness = 1pt,
      margin-bottom = 10pt,
      artifact-type = layout,
    }
    \decorateltxtalkarea{footer}{
      frame = top,
      rule-color = ltxtalk-primary!30,
      rule-thickness = 0.4pt,
      margin-top = 5pt,
      artifact-type = decorative,
    }
    
    % === FARBEN ===
    \setltxtalkcolors{
      primary   = {HTML}{333333},
      secondary = {HTML}{666666},
      accent    = {HTML}{999999},
      text      = {HTML}{222222},
      background = {HTML}{FFFFFF},
    }
  }

\RegisterLTXTalkTheme{minimal}{\__ltxtalk_minimal_theme_setup:}

%</minimal>
%<*magpie>
% \subsection{Magpie theme (inspired by Beamer Madrid)}
\NeedsTeXFormat{LaTeX2e}
\ProvidesExplPackage{ltxtalk-theme-magpie}
{2026/08/05}
{v0.4.1}
{Magpie theme for ltx-talk, inspired by Beamer Madrid}
\RequirePackage{ltxtalk-theme}
\cs_new_protected:Npn \__ltxtalk_magpie_theme_setup:
  {
    \ltxtalksetup{header=true,header-height=1.15cm,footer=true,footer-height=0.48cm}
    \setltxtalktitleplacement{frametitle=header,section=none}
    \defineltxtalkslot{magpie-frametitle}{header}{frametitle}
      {role=heading,heading-level=2}
    \defineltxtalkslot{magpie-framesubtitle}{header}{framesubtitle}
      {role=heading,heading-level=3}
    \defineltxtalkslot{magpie-author}{footer}{custom}{role=note}
    \defineltxtalkslot{magpie-title}{footer}{custom}{role=note}
    \defineltxtalkslot{magpie-page}{footer}{custom}{role=navigation}
    \setltxtalkslotcontent{magpie-author}{\@shortauthor\ (\@shortinstitute)}
    \setltxtalkslotcontent{magpie-title}{\@shorttitle}
    \setltxtalkslotcontent{magpie-page}{\@shortdate\quad\@framenumber}
    \positionltxtalkslot{magpie-frametitle}{0.22cm}{0.22cm}
      {\dim_eval:n{\paperwidth-0.44cm}}{0.62cm}
    \positionltxtalkslot{magpie-framesubtitle}{0.22cm}{0.72cm}
      {\dim_eval:n{\paperwidth-0.44cm}}{0.38cm}
    \positionltxtalkslot{magpie-author}{0.18cm}{center}{0.30\paperwidth}{0.42cm}
    \positionltxtalkslot{magpie-title}{center}{center}{0.34\paperwidth}{0.42cm}
    \positionltxtalkslot{magpie-page}
      {\dim_eval:n{0.88\paperwidth-0.18cm}}{center}{0.12\paperwidth}{0.42cm}
    \styleltxtalkslot{magpie-frametitle}{font=\Large\bfseries\sffamily,color=white,
      align=left,padding=3pt}
    \styleltxtalkslot{magpie-framesubtitle}
      {font=\small\sffamily,color=white,align=left}
    \styleltxtalkslot{magpie-author}{font=\tiny\sffamily,color=white,align=left,
      background=ltxtalk-secondary}
    \styleltxtalkslot{magpie-title}{font=\tiny\sffamily,color=white,align=center,
      background=ltxtalk-primary!85!black}
    \styleltxtalkslot{magpie-page}{font=\tiny\sffamily,color=white,align=right}
    \decorateltxtalkarea{header}{background=ltxtalk-primary,artifact-type=layout}
    \decorateltxtalkarea{footer}{background=ltxtalk-primary!70!black,artifact-type=layout}
    \AssignSocketPlug { lttheme/page/background } { magpie-title }
    \AssignSocketPlug { lttheme/page/foreground } { magpie-title }
    \EditTemplateDefaults {titlepage}{talk}
      {framestyle=lttheme,vertical-alignment=top}
    \EditInstance {titlepage-element}{title}
      {color=white,font=\Large\sffamily,before-skip=1.35cm,after-skip=0.1em}
    \EditInstance {titlepage-element}{subtitle}
      {color=white,font=\small\sffamily,before-skip=0pt,after-skip=0.8em}
    \EditInstance {titlepage-element}{author}
      {color=ltxtalk-text,before-skip=0.8em,after-skip=0.7em}
    \EditInstance {titlepage-element}{institute}
      {color=ltxtalk-text,font=\scriptsize,after-skip=0.7em}
    \EditInstance {titlepage-element}{date}{color=ltxtalk-text}
    \setltxtalkcolors{primary={rgb}{0.2,0.2,0.7},secondary={rgb}{0.15,0.15,0.525},
      accent={rgb}{0.2,0.2,0.7},text={HTML}{000000},background={HTML}{FFFFFF}}
  }
\RegisterLTXTalkTheme{magpie}{\__ltxtalk_magpie_theme_setup:}
%</magpie>
%<*hawk>
% \subsection{Hawk theme (inspired by Beamer Hannover)}
\NeedsTeXFormat{LaTeX2e}
\ProvidesExplPackage{ltxtalk-theme-hawk}
{2026/08/05}
{v0.4.1}
{Hawk theme for ltx-talk, inspired by Beamer Hannover}
\RequirePackage{ltxtalk-theme}
\cs_new_protected:Npn \__ltxtalk_hawk_theme_setup:
  {
    \ltxtalksetup{header=true,header-height=1.05cm,left-margin=true,
      margin-width=1.6cm,footer=false}
    \setltxtalktitleplacement{frametitle=header,section=none}
    \defineltxtalkslot{hawk-title}{left}{custom}{role=note}
    \defineltxtalkslot{hawk-author}{left}{custom}{role=note}
    \defineltxtalkslot{hawk-section}{left}{section}{role=navigation}
    \defineltxtalkslot{hawk-subsection}{left}{subsection}{role=navigation}
    \defineltxtalkslot{hawk-frametitle}{header}{frametitle}{role=heading,heading-level=2}
    \defineltxtalkslot{hawk-framesubtitle}{header}{framesubtitle}{role=heading,heading-level=3}
    \setltxtalkslotcontent{hawk-title}{\@shorttitle}
    \setltxtalkslotcontent{hawk-author}{\@shortauthor}
    \positionltxtalkslot{hawk-title}{center}{0.25cm}{1.45cm}{0.55cm}
    \positionltxtalkslot{hawk-author}{center}{0.9cm}{1.45cm}{0.5cm}
    \positionltxtalkslot{hawk-section}{0.12cm}{1.8cm}{1.36cm}{0.7cm}
    \positionltxtalkslot{hawk-subsection}{0.12cm}{2.65cm}{1.36cm}{0.7cm}
    \positionltxtalkslot{hawk-frametitle}
      {\dim_eval:n{0.58\paperwidth-2cm}}{0.12cm}{0.42\paperwidth}{0.72cm}
    \positionltxtalkslot{hawk-framesubtitle}
      {\dim_eval:n{0.58\paperwidth-2cm}}{0.78cm}{0.42\paperwidth}{0.42cm}
    \styleltxtalkslot{hawk-title}{font=\fontsize{5}{6}\selectfont\sffamily,
      color=ltxtalk-text,align=center}
    \styleltxtalkslot{hawk-author}{font=\tiny\sffamily,color=ltxtalk-secondary,align=center}
    \styleltxtalkslot{hawk-section}{font=\small\bfseries\sffamily,color=ltxtalk-primary,align=left}
    \styleltxtalkslot{hawk-subsection}{font=\scriptsize\sffamily,color=ltxtalk-secondary,align=left}
    \styleltxtalkslot{hawk-frametitle}{font=\Large\sffamily,color=ltxtalk-primary,align=right}
    \styleltxtalkslot{hawk-framesubtitle}{font=\small\sffamily,color=ltxtalk-secondary,align=right}
    \decorateltxtalkarea{left}{background=ltxtalk-primary!20,artifact-type=layout}
    \AssignSocketPlug { lttheme/area/left } { hawk }
    \EditTemplateDefaults {titlepage}{talk}{framestyle=lttheme,vertical-alignment=center}
    \EditInstance {titlepage-element}{title}
      {color=ltxtalk-primary,font=\Large\sffamily,after-skip=0.15em}
    \EditInstance {titlepage-element}{subtitle}
      {color=ltxtalk-primary,font=\small\sffamily,before-skip=0pt,after-skip=0.9em}
    \EditInstance {titlepage-element}{author}{color=ltxtalk-text,before-skip=0.8em}
    \EditInstance {titlepage-element}{institute}{color=ltxtalk-text,font=\scriptsize}
    \EditInstance {titlepage-element}{date}{color=ltxtalk-text}
    \setltxtalkcolors{primary={rgb}{0.2,0.2,0.7},secondary={rgb}{0.15,0.15,0.525},
      accent={rgb}{0.2,0.2,0.7},text={HTML}{000000},background={HTML}{FFFFFF}}
  }
\RegisterLTXTalkTheme{hawk}{\__ltxtalk_hawk_theme_setup:}
%</hawk>
%<hawk>\endinput

% \subsection{Sparrow theme (inspired by Beamer Szeged)}
%<*sparrow>
\NeedsTeXFormat{LaTeX2e}
\ProvidesExplPackage{ltxtalk-theme-sparrow}
{2026/08/05}
{v0.4.1}
{Sparrow theme for ltx-talk, inspired by Beamer Szeged}
\RequirePackage{ltxtalk-theme}
\cs_new_protected:Npn \__ltxtalk_sparrow_theme_setup:
  {
    \ltxtalksetup{header=true,header-height=1.35cm,footer=true,footer-height=0.42cm}
    \setltxtalktitleplacement{frametitle=content,section=header}
    \defineltxtalkslot{sparrow-section}{header}{section}{role=navigation}
    \defineltxtalkslot{sparrow-subsection}{header}{subsection}{role=navigation}
    \defineltxtalkslot{sparrow-frametitle}{header}{frametitle}{role=heading,heading-level=2}
    \defineltxtalkslot{sparrow-framesubtitle}{header}{framesubtitle}{role=heading,heading-level=3}
    \defineltxtalkslot{sparrow-institute}{footer}{institute}{role=note}
    \defineltxtalkslot{sparrow-title}{footer}{title}{role=note}
    \defineltxtalkslot{sparrow-page}{footer}{pagenumber}{role=navigation}
    \positionltxtalkslot{sparrow-section}{0.18cm}{0.02cm}{0.44\paperwidth}{0.3cm}
    \positionltxtalkslot{sparrow-subsection}
      {\dim_eval:n{0.56\paperwidth-0.18cm}}{0.02cm}{0.44\paperwidth}{0.3cm}
    \positionltxtalkslot{sparrow-frametitle}{0.25cm}{1.38cm}
      {\dim_eval:n{\paperwidth-0.5cm}}{0.58cm}
    \positionltxtalkslot{sparrow-framesubtitle}{0.25cm}{1.84cm}
      {\dim_eval:n{\paperwidth-0.5cm}}{0.35cm}
    \positionltxtalkslot{sparrow-institute}{0.18cm}{center}{0.3\paperwidth}{0.45cm}
    \positionltxtalkslot{sparrow-title}{center}{center}{0.4\paperwidth}{0.45cm}
    \positionltxtalkslot{sparrow-page}
      {\dim_eval:n{0.88\paperwidth-0.18cm}}{center}{0.12\paperwidth}{0.45cm}
    \styleltxtalkslot{sparrow-section}{font=\scriptsize\bfseries\sffamily,color=ltxtalk-primary,align=left}
    \styleltxtalkslot{sparrow-subsection}{font=\scriptsize\sffamily,color=ltxtalk-secondary,align=right}
    \styleltxtalkslot{sparrow-frametitle}
      {font=\Large\sffamily,color=ltxtalk-primary,align=left}
    \styleltxtalkslot{sparrow-framesubtitle}
      {font=\small\sffamily,color=ltxtalk-secondary,align=left}
    \styleltxtalkslot{sparrow-institute}{font=\tiny\sffamily,color=ltxtalk-secondary,align=left}
    \styleltxtalkslot{sparrow-title}{font=\tiny\sffamily,color=ltxtalk-primary,align=center}
    \styleltxtalkslot{sparrow-page}{font=\tiny\sffamily,color=ltxtalk-secondary,align=right}
    \decorateltxtalkarea{header}{background=white,artifact-type=layout}
    \decorateltxtalkarea{footer}{frame=top,rule-color=ltxtalk-primary!50,
      rule-thickness=1pt,background=ltxtalk-primary!5,artifact-type=layout}
    \AssignSocketPlug { lttheme/page/foreground } { sparrow-bars }
    \EditTemplateDefaults {titlepage}{talk}{framestyle=lttheme,vertical-alignment=center}
    \EditInstance {titlepage-element}{title}
      {color=ltxtalk-primary,font=\Large\sffamily,after-skip=0.15em}
    \EditInstance {titlepage-element}{subtitle}
      {color=ltxtalk-primary,font=\small\sffamily,before-skip=0pt,after-skip=0.9em}
    \EditInstance {titlepage-element}{author}{color=ltxtalk-text,before-skip=0.8em}
    \EditInstance {titlepage-element}{institute}{color=ltxtalk-text,font=\scriptsize}
    \EditInstance {titlepage-element}{date}{color=ltxtalk-text}
    \setltxtalkcolors{primary={rgb}{0.2,0.2,0.7},secondary={rgb}{0.15,0.15,0.525},
      accent={rgb}{0.2,0.2,0.7},text={HTML}{000000},background={HTML}{FFFFFF}}
  }
\RegisterLTXTalkTheme{sparrow}{\__ltxtalk_sparrow_theme_setup:}
%</sparrow>
%<*anchovy>
% \subsection{Anchovy colour theme (colours inspired by Beamer AnnArbor)}
\NeedsTeXFormat{LaTeX2e}
\ProvidesExplPackage{ltxtalk-theme-anchovy}
{2026/08/05}
{v0.4.1}
{Anchovy colour theme for ltx-talk, inspired by Beamer AnnArbor}
\RequirePackage{ltxtalk-theme}
\cs_new_protected:Npn \__ltxtalk_anchovy_theme_setup:
  {\setltxtalkcolors{primary={rgb}{0,0,0.8},secondary={HTML}{65500A},
    accent={HTML}{B45309},text={HTML}{111827},background={HTML}{FFF7D6}}}
\RegisterLTXTalkTheme{anchovy}{\__ltxtalk_anchovy_theme_setup:}
%</anchovy>
%<*carp>
% \subsection{Carp colour theme (colours inspired by Beamer CambridgeUS)}
\NeedsTeXFormat{LaTeX2e}
\ProvidesExplPackage{ltxtalk-theme-carp}
{2026/08/05}
{v0.4.1}
{Carp colour theme for ltx-talk, inspired by Beamer CambridgeUS}
\RequirePackage{ltxtalk-theme}
\cs_new_protected:Npn \__ltxtalk_carp_theme_setup:
  {\setltxtalkcolors{primary={rgb}{0.8,0,0},secondary={HTML}{7F1D1D},
    accent={HTML}{991B1B},text={HTML}{111111},background={HTML}{FFFFFF}}}
\RegisterLTXTalkTheme{carp}{\__ltxtalk_carp_theme_setup:}
%</carp>
%<carp>\endinput
